

\section{Introduction}
\label{sec:intro}

Given a bipartite graph $G(U,V;E)$, where $U$ and $V$ are the sets of vertices and $E \subset U \times V$ is the set of edges, a matching $M \subset E$ is a set of edges such that every vertex is incident with at most one edge of $M$. For simplicity of notation, for every $n$ we shall only consider the following class of bipartite graphs, that we shall refer to as $G_n$. For every $G(U,V;E) \in G_n$ it holds that $|U| = |V| = n$ and that $E$ contains a matching of size $n$ (and hence $G$ has a perfect matching). All results that we will state for $G_n$ hold without change for all bipartite graphs that have a matching of size $n$ (and arbitrarily many vertices).


Karp, Vazirani and Vazirani~\cite{KarpVV1990} introduced the \emph{online bipartite matching} problem. In this problem, a bipartite graph $G(U,V;E)$ is revealed in an online manner, where side $V$ is known in advance, and vertices in $U$ are revealed one after another, along with their edges to $V$. When a vertex $u \in U$ is revealed together with its set of neighbors $N(u) \subset V$, the algorithm must decide, immediately and irrevocably, which of the yet unmatched vertices in $N(u)$ to match to $u$ (if any). The goal of the algorithm is to maximize the size of the matching achieved.

This setting can be viewed as a game between two players: a maximizing player who wishes the resulting matching to be as large as possible, and a minimizing player who wishes the matching to be as small as possible. First, the minimizing player chooses $G(U,V;E)$ in private (without the maximizing player seeing $E$), subject to $G \in G_n$. Thereafter, the structure of $G$ is revealed to the maximizing player in $n$ steps, where at step $j$ (for $1 \le j \le n$) the set $N(u_j) \subset V$ of vertices adjacent to $u_j$ is revealed. At every step $j$, upon seeing $N(u_j)$ (and based on all edges previously seen and all previous matching decisions made), the maximizing player needs to irrevocably either match $u_j$ to a currently unmatched vertex in $N(u_j)$, or leave $u_j$ unmatched.

There %
has been   %
much recent interest in the online bipartite matching problem and variations and generalizations of it, as such models have applications 
%
to     %
allocation problems in certain economic settings, in which buyers (vertices of $U$) arrive online and are interested in purchasing various items (vertices of $V$). A prominent example of such an application is online advertising; for more details, see for example the survey by Mehta~\cite{Mehta13}. The new problems are both theoretically elegant and practically relevant.

\paragraph*{Max-min greedy matching} We study a setting related to online bipartite matching, that we call
%
%
%
%
\emph{max-min greedy matching}.   %
Our setting is also a game between a maximizing player and a minimizing player. The bipartite graph $G(U,V;E)\in G_n$ is given upfront. Upon seeing $G$ the maximizing player chooses a linear order $\pi$ over $V$. Upon seeing $G$ and $\pi$, the minimizing player chooses a linear order $\sigma$ over $U$.
The combination of $G$, $\pi$ and $\sigma$ %
defines   %
a unique matching $M_G[\sigma,\pi]$ that we refer to as the greedy matching.
It is the matching produced by the greedy matching algorithm in which vertices of $U$ arrive in order $\sigma$ and each vertex $u \in U$ is matched to the highest (under $\pi$) yet unmatched $v \in N(u)$ (or left unmatched, if
%
all of $N(u)$ has already been matched).   %

The matching $M_G[\sigma,\pi]$ has several additional equivalent definitions.
For example, $M_G[\sigma,\pi]$ is the matching produced by the greedy matching algorithm in which vertices of $V$ arrive in order $\pi$ and each vertex $v \in V$ is matched to the highest (earliest in arrival order under $\sigma$) yet unmatched $u \in N(v)$ (or left unmatched, if
%
all of $N(v)$ has already been matched).  %
Also, $M_G[\sigma,\pi]$ is the unique stable matching in $G$ (in the sense of~\cite{GS62}), if the preference order of every vertex $u \in U$ over its neighbors is consistent with $\pi$, and the preference order of every vertex $v \in V$ over its neighbors is consistent with $\sigma$. More details are given in \expref{Section}{sec:stable-matching}.

Let $$\rho[G] = \frac{1}{n} \max_{\pi} \min_{\sigma}[|M_G[\sigma,\pi]|],$$ and let $$\rho = \min_{G \in G_n}[\rho[G]].$$
It is easy to see that $\rho \ge 1/2$.
In fact, to ensure a matching of size $n/2$, the max player need not work hard.
Since every greedy matching is a \emph{maximal} matching, for \emph{every} linear order $\pi$ the obtained matching is of size at least $n/2$.
The question we study %
here  %
is whether the max player can ensure a matching of size strictly 
greater than $n/2$; that is, whether $\rho$ is bounded away from $1/2$.  
%
%
%

For an upper bound on $\rho$, it was observed by Cohen-Addad et al.~\cite{Cohen-AddadEFF16} that $\rho \leq 2/3$.
To show this, they observe that in the 6-cycle graph, depicted in \expref{Figure}{fig:two-thirds}, no linear order $\pi$ can guarantee to match more than two vertices in the worst case. Indeed, suppose (without loss of generality) that $\pi=(v_3,v_2,v_1)$. For $\sigma=(u_1,u_3,u_2)$, $u_1$ is matched to $v_2$, $u_3$ is matched to $v_3$, and $u_2$ is left unmatched, resulting in a matching of size $2$.

\begin{figure}[h!]
	\begin{center}
\includegraphics[trim={0 7.5cm 0 4cm},clip,scale=.5]{figure1}
		%
		\caption{For every linear order $\pi$ there exists a linear order $\sigma$ that matches only 2 of the 3 vertices. Thick edges are in the matching; gray vertices are unmatched.}
		\label{fig:two-thirds}
	\end{center}
\end{figure}

%

\subsection{Our results}
\label{sec:results}

Our main result resolves the open problem in the affirmative. 
%

\vspace{0.1in}
\noindent {\bf Main result [\expref{Theorem}{thm:86}]:}
%
%
\emph{For all graphs,  %
$\rho \geq 1/2 + 1/86 > 0.51$.
Moreover, there is a polynomial-time algorithm that given $G(U,V;E)$
produces a linear order $\pi$ over $V$ satisfying the above bound.}
\vspace{0.1in}

The significance of this result is that $1/2$ is not the optimal answer.
We believe that further improvements are possible.
In fact, for Hamiltonian graphs we show that $\rho \ge 5/9$ (see \expref{Section}{sec:hamiltonian}).

The proof method is quite involved; it is natural to ask whether simpler approaches may work.
In what follows we specify three natural attempts that all fail.

\paragraph*{Failed attempt 1: random linear order.}
A first attempt would be to check whether a random linear order $\pi$ obtains the desired result (in expectation)\footnote{Unless stated otherwise, whenever way say ``random linear order,'' we refer to choosing a linear order uniformly at random.}.
The performance of a random linear order is interesting for an additional reason: it is the performance in scenarios where the graph structure is unknown to the designer.
Unfortunately, there exists a bipartite graph $G$, even one where all vertices have high degree, for which a random linear order matches no more than a
%
$1/2 + o(1)$ fraction   %
of the vertices (see \expref{Section}{sec:random}).

In contrast, we show that in the case of Hamiltonian graphs a random linear order guarantees a competitive ratio strictly greater than $1/2$ (\expref{Section}{sec:random}).
A similar proof approach applies to regular graphs as well.

\vspace{0.1in}
%
\noindent {\bf Results for a random linear order [\expref{Theorem}{thm:random-hamiltonian}]:}   %
%
\emph{There is some constant $\rho_0^{\mathrm{H}} > 1/2$ such that for every Hamiltonian graph $G \in G_n$, regardless of $n$,
a random linear order $\pi$ results in $\rho \geq \rho_0^{\mathrm{H}}$.
Similarly, there is some constant $\rho_0^{\mathrm{R}} > 1/2$ such that for every $d$-regular graph $G$, regardless of $d,n$, a random linear order $\pi$ results in $\rho \geq \rho_0^{\mathrm{R}}$.}
\vspace{0.1in}

\paragraph*{Failed attempt 2: iterative upgrading.}
A second attempt would be to iteratively ``upgrade'' unmatched vertices, with the hope that the iterative process will reach a state where many vertices will be matched.
That is, in every iteration consider the worst order $\sigma$ for the current linear order $\pi$ and move all unmatched vertices (in the matching
%
produced %
by $(\pi,\sigma)$) to be ranked higher in $\pi$. This algorithm is similar to the \emph{$k$-pass Category Advice} algorithm of~\cite{Bor2017}, but with the difference that in~\cite{Bor2017} $\sigma$ remains unchanged throughout the $k$ iterations. In~\cite{Bor2017} it was shown that in their setting, the fraction of matched vertices approaches $2/(1 + \sqrt{5}) \simeq 0.619$ as $k$ grows. In contrast,
in \expref{Section}{sec:iterative} we show that in our setting this process can go on for $\log n$ iterations before reaching a linear order that matches more than a half of the vertices.
This fact gives some indication that establishing a proof using this operator might be difficult.

\paragraph*{Failed attempt 3: degree-based ranking.}
A third attempt would be to give preference to vertices with lower degrees, as they would have fewer opportunities to be matched to incoming vertices of $U$. %
Consider a graph with multiple copies of the subgraph $(u_1,v_1),(u_2,v_2),(u_1,v_2)$ along with two additional vertices $u_a,u_b$ (and their partners $v_a,v_b$). If we connect all vertices of type $v_1$ to $u_a$ and $u_b$, we get that their degree is $3$, while the degree of vertices of type $v_2$ is $2$. If $\pi$ is chosen according to the degree, vertices of type $v_2$ will be ranked higher than vertices of type $v_1$. In this case, if $\sigma$ orders the vertices of type $u_1$ first, they will be matched to vertices of type $v_2$, leaving the vertices of type $v_1$ unmatched. The resulting matching will therefore be of size $(1/2+o(1))n$.


%
%
%
%
%
%

%

\paragraph*{Why is this model interesting mathematically?} The setting of max-min greedy matching is easy to state.
%
%
%
The deceptively simple problem of getting a ratio bounded away
from one half turns out to be quite difficult.   %

As discussed above, several natural approaches fail to achieve this. The problem remained open for quite some time, despite attempts to solve it. Indeed, the solution that we find is not simple; it involves taking the best of four algorithms. However, these algorithms are not unrelated. They all share a unifying theme that involves a clean combinatorial property, referred to as a maximal path cover (see \expref{Section}{sec:main-result}). This theme enabled us to break the barrier of half, but interesting problems remain open, such as whether the bound of 2/3 can be achieved. We hope that the progress made in this
%
paper   %
will motivate and enable further improvement in this interesting problem.

\subsection{Additional results}
We further establish lower and upper bounds for regular graphs.

\vspace{0.1in}
\noindent {\bf Results for regular graphs [\expref{Corollary}{cor:regular} and \expref{Theorem}{thm:reg_89}]:}
%
\emph{For $d$-regulars bipartite graphs, $\rho \ge 5/9 - O(1/\sqrt{d})$.
On the other hand, for every integer $d \ge 1$, there is a regular graph $G_{d}$ of even degree $2d$ such that $\rho(G_{d}) \le 8/9$.}
\vspace{0.1in}

An additional natural problem is to find the best linear order $\pi$, given a graph $G$.
We suspect that this is a difficult computational problem.
However, the special case of determining whether there is a \emph{perfect} $\pi$ (a linear order on $V$ that for every linear order $\sigma$ leads to a perfect matching) does have a polynomial-time algorithm.  %
%
%
(The proof appears in \expref{Section}{sec:perfect-pi}.)

\vspace{0.1in}
\begin{theorem}
	\label{thm:perfect-pi}
There is a polynomial-time algorithm that given a graph $G \in G_n$ determines whether $G$ has a \emph{perfect} $\pi$, and if so,
%
computes   %
a perfect $\pi$.
\end{theorem}

%
%

\subsection{Application to resource allocation and pricing}
\label{sec:interpretations}

Various problems related to online bipartite matching are closely related to problems that attract attention in the algorithms community. Gaining better understanding of the max-min greedy matching problem sheds light on more general problems, some of which are still open. In what follows we elaborate on an application to a pricing problem.

Feldman et al. \cite{FeldmanGL15} study the design of pricing mechanisms for allocation of items in markets.
The basic setting is a matching market, where $v_{ij}$ is the
%
value agent $i$ obtains from getting %
item $j$, and every agent can receive at most a single item. The seller assigns prices to items, and agents arrive in an adversarial order (after observing the prices), each purchasing an (arbitrary) item that maximizes their utility (defined as value minus price).
It is shown that, given a weighted bipartite graph (with agents on one side, items of the other side, and weight $v_{ij}$ for the edge between agent $i$ and item $j$), one can set item prices that guarantee at least half of the optimal welfare for any arrival order.
The last result holds in much more general settings, namely settings where buyers have submodular valuations over bundles of items~\footnote{A valuation is said to be submodular if for every two sets $S,T$, $v(S)+v(T) \geq v(S\cup T) + v(S \cap T)$,.}, and even in a Bayesain settings, where the seller knows only the (product) distribution from which agent valuations are drawn, but not their realizations.

In the Bayesian setting no item prices can guarantee better than half of the optimal welfare in the worst case.
A natural question is whether this ratio can be improved in scenarios where the designer knows the realized values of the buyers from the outset~\footnote{The full information assumption is sensible in repeated markets or in markets where the stakes are high and the designer may invest in learning the demand in the market before setting prices.}.
Concretely, do there exist item prices that guarantee strictly more than half the optimal welfare, for any arrival order $\sigma$?
Not only has this question been open for general combinatorial auctions with submodular valuations, it has been open even for unit-demand buyers, and even if all individual values are in $\{0,1\}$ (henceforth referred to as binary unit demand valuations).
In the latter setting, pricing is equivalent to imposing a linear order over the items, hence the max-min greedy matching is a precise formulation for the pricing problem in binary unit-demand settings.

An equivalent scenario is one where in each step the ``items player'' offers an item, and the ``buyers player,'' %
upon seeing the item, allocates the item to one of the buyers who
%
want     %
the item (if there %
are %
any), and that buyer leaves.
The items player is non-adaptive (plays blindfolded, without seeing which buyers remain\footnote{We note that when the items player is adaptive (chooses the next item based on what happened in the past), the items player can ensure a perfect matching. This is done as follows: in each step, find a minimal tight set of items, and offer an arbitrary item from that set. Here, a set of items is tight if the number of buyers that want items in the set is equal to the size of the set.}
).
The size of the matching that can be guaranteed by the items player is equivalent to the max-min greedy matching problem.

Yet an additional equivalent formulation of the problem is one where the linear order $\pi$ is imposed over the buyers rather than over the items.
The buyers then arrive in the order of $\pi$, each taking an arbitrary item
%
they want.    %
One can verify that the size of the matching that can be guaranteed by an ordering over the buyers is equivalent to the max-min greedy matching problem.


\subsection{Relation to prior work}
\label{sec:related-work}

\paragraph*{Relation to online bipartite matching}

The max-min greedy Matching problem %
%
%
%
is a nonstandard version of online problems. In the standard online matching problem~\cite{KarpVV1990}, the algorithm designer has control over the matching algorithm, but has no control over the arrival order of clients (vertices). Our setting can model a situation in which the designer (the maximizing player) has full control over the arrival order of clients (it knows which ``items'' in $U$ each ``client'' in $V$ wants, and it chooses $\pi$ based on this knowledge), but no control over the matching algorithm (the minimizing player can choose the worst possible match in every step, effectively resulting in a linear order $\sigma$).

Karp et al. \cite{KarpVV1990} introduced the \emph{Ranking} algorithm which
has a %
$1-1/\eee$   %
competitive ratio in the online bipartite matching setting \cite{KarpVV1990,Goel2008}. Translating this algorithm to our max-min greedy setting, it amounts to simply selecting $\pi$ at random, and then the minimizing player selects $\sigma$ after seeing $\pi$. We show that there are bipartite graphs $G \in G_n$ for which with high probability over the random choice of $\pi$, there is a choice of $\sigma$ resulting in $M_G[\sigma,\pi] \le 1/2 + o(1)$. Karp et al. \cite{KarpVV1990} also showed that no algorithm for online bipartite matching has a competitive ratio better than %
$1 - 1/\eee + o(1)$.   %
This was shown by exhibiting a distribution over ``difficult'' graphs. Each graph in the support of this distribution has a unique perfect matching, and consequently (see \expref{Proposition}{prop:dag}), there is a linear order $\pi$ in the max-min greedy setting that ensures that all vertices are matched (regardless of $\sigma$). Hence neither the lower bounds nor the upper bounds known for the online matching model give useful bounds in the max-min greedy model.

There are additional known results for online bipartite matching.
For $d$-regular graphs, Cohen and Wajc \cite{CohenW18} present a
%
randomized  %
algorithm that obtains $1-O(\sqrt{\log d}/\sqrt{d})$ in expectation, and a lower bound of $1-O(1/\sqrt{d})$. This is in contrast to our \expref{Theorem}{thm:reg_89} that shows that $\rho$ is bounded away from~$1$ even when $d$ is arbitrarily large.
For general bipartite graphs, under random (rather than adversarial) arrival order, the deterministic greedy algorithm gives %
$1-1/\eee$    %
and no deterministic algorithm can obtain more than $3/4$ \cite{Goel2008}.
\emph{Ranking} (which is a randomized algorithm) obtains at least $0.696$ of the optimal matching \cite{Mahdian2011} and at most $0.727$ \cite{Karande2011}.
No %
randomized   %
algorithm can obtain more than $0.823$ \cite{Manshadi2012}.

\paragraph*{Relation to pricing mechanisms}
Our work is also related to the recent body of literature on pricing mechanisms.
Motivated by the fact that in real-life situations one is often willing to trade optimality for simplicity, the study of simple mechanisms has gained a lot of interest in the literature on algorithmic mechanism design.
One of the simplest forms of mechanisms is that of posted price mechanisms, where prices are associated with items and agents buy their most preferred bundles as they arrive.
Pricing mechanisms have many advantages: they are simple, straightforward, and allow for asynchronous arrival and departure of buyers. %
Various forms of posted price mechanisms for welfare maximization have been proposed for various combinatorial settings \cite{FeldmanGL15,DuttingFKL16,Lucier17,EzraFRS20}. These mechanisms are divided along several axes, such as item vs. bundle pricing, static vs. dynamic pricing, and anonymous vs. personalized pricing.
For any market with submodular valuations, one can obtain $1/2$ of the optimal welfare by static, anonymous item prices \cite{FeldmanGL15}.
Until the present paper, no %
results better than  %
$1/2$ were known even for markets with unit-demand valuations with $\{0,1\}$ individual values.
For a market with $m$ identical items, there exists a pricing scheme that obtains at least $5/7-1/m$ of the optimal welfare for submodular valuations \cite{EzraFRS20}.






\section{Proof of main result}
\label{sec:main-result}

The graph $G(U,V;E)$ with $|U| = |V| = n$ has a perfect matching $M$ in which $u_i \in U$ is matched with $v_i \in V$ for every $1 \le i \le n$. For a given $i$, we refer to $u_i$ and $v_i$ as \emph{partners} of each other. Given a set $S \subset V$, the set of neighbors of $S$ is denoted  as $N(S)$ (where necessarily $N(S) \subset U$).
%
In this section we prove our main result.

\begin{theorem}
	\label{thm:86}
	Given a bipartite graph $G(U,V;E)$ with a perfect matching $\{(u_i,v_i)\}$, there exists a linear order $\pi$ that guarantees that the greedy matching will be of size at least $\frac{22}{43}n$, regardless of $\sigma$. Moreover, there is a polynomial-time algorithm that chooses $\pi$ with such a guarantee.
\end{theorem}

Our proof approach is as follows. We shall first associate with $G$ an auxiliary directed graph that we refer to as the 
\emph{spoiling graph} %
%
$H(V,D)$. This notion by itself is not
%
new---similar notions   %
%
appear   %
in previous related work. The new aspect related to the spoiling graph and the key to our approach is a notion of a \emph{maximal path cover}. Given a maximal path cover of the spoiling graph (which, %
as we show in \expref{Proposition}{pro:maximalcover}, can be found in polynomial time), we partition the set $V$ of vertices into four classes, depending on their roles in the maximal path cover. The classes are $V_1$ (singleton vertices), $S$ (start vertices of paths), $T$ (end vertices of paths), and $I$ (intermediate vertices of paths). By considering several carefully chosen orders among the classes of vertices, and also of vertices within the classes, we obtain four possible candidate linear orders for $\pi$, denoted $\pi_1, \pi_2, \pi_3, \pi_4$. We show that for every bipartite graph with a perfect matching, at least one of these linear orders, if used as $\pi$, guarantees that the greedy matching will be of size at least $\frac{22}{43}n$, for every $\sigma$.
Put in other words, if for each of $\{\pi_1, \pi_2, \pi_3, \pi_4\}$ there is a linear order over $U$ for which the greedy matching is smaller than $\frac{22}{43}n$, this would imply (using properties listed in \expref{Lemma}{lem:properties}) that the path cover giving rise to these linear orders was not maximal. Our lemma statements provide tools to identify the correct $\pi$ in polynomial time using properties of the classes $V_1, S, T,$ and $I$ described above.


\begin{figure}[h!]
	\begin{center}
\includegraphics[trim={0 3.5cm 0 4cm},clip,scale=.5]{figure2}		
\caption{(a) A graph $G(U,V;E)$; (b) The corresponding spoiling graph of $G$, $H(V,D)$.}
	\label{fig:spoiling}
	\end{center}
\end{figure}


We now proceed to define the spoiling graph. Given $G(U,V;E)$,
consider a directed graph $H(V,D)$ whose vertices are the set $V$, and whose set $D$ of directed edges (arcs) is defined as follows: $(v_i,v_j) \in D$ iff $(u_j,v_i) \in E$ (see \expref{Figure}{fig:spoiling} for an illustration). %
We refer to $H(V,D)$ as the \emph{spoiling graph} for $G$, because arc $(v_i,v_j) \in D$ allows for the possibility that edge $(u_j,v_i)\in E$ is chosen into a matching $M'$ in $G$, spoiling for $v_j$ the possibility (offered by $M$) of being matched to $u_j$. Note that this spoiling effect may materialize in a $(\sigma,\pi)$ matching only if $v_i$ is ranked higher than $v_j$ in $\pi$. Hence the spoiling graph conveys information that may be relevant to the choice of $\pi$.
%
%

As an example of the information that can be derived from the spoiling graph, consider the following proposition (whose proof can be also obtained as a special case of a result given in \cite{Cohen-AddadEFF16} and \cite{LemeW17} for the more general case of Gross Substitutes valuations).

\begin{proposition}
	\label{prop:dag}
	If $G$ has a unique perfect matching, then $\rho(G)=1$.
\end{proposition}
\begin{proof}
	Let $u_i\in U$ and $v_i \in V$ be partners in the unique perfect matching $M$ of $G$.
	We claim that the spoiling graph $H$ of $G$ is a \emph{directed acyclic graph} (DAG).
	Suppose toward a %
contradiction that $H$ contains a simple directed cycle $v_{i_1},v_{i_2},\ldots, v_{i_\ell},v_{i_1}$.
	This directed cycle corresponds to the cycle $u_{i_1},v_{i_1},u_{i_2},v_{i_2},\ldots,u_{i_\ell},v_{i_{\ell}},u_{i_1}$ in $G$.
	But removing the edges $(u_{i_j},v_{i_j})$, $1\leq j\leq \ell$ from $M$ and adding the edges $(u_{i_j},v_{i_{j+1}})$ to $M$ (where $\ell+1=1$) yields a different perfect matching, contradicting the uniqueness of $M$.
	
	Since $H$ is a DAG, we can topologically sort its vertices and choose a linear order $\pi$ such that earlier vertices in the topological order have a lower rank in $\pi$. This ensures that for every directed edge $(v,w)$ in $H$,
%
the partner of vertex $v$  %
will never prefer $w$ over $v$. Thus, every vertex chooses its partner in $M$ upon arrival.
\end{proof}

We now proceed to define the notion of a maximal path cover.
A \emph{directed path} $P$ (whose length is denoted by $|P|$) in $H$ is a sequence of $|P|$ vertices (say, $v_1, \ldots, v_{|P|}$) such that $(v_{i}, v_{i+1}) \in D$ for all $1 \le i \le |P|-1$.  
A single vertex is a path of length~1. %
%
%
%
%
%
%
%
A \emph{path cover} of $H$ is a collection of
%
vertex-disjoint   %
directed paths that covers all vertices in $V$. We consider the following operations that can transform a given path cover to a different one (see \expref{Figure}{fig:path-cover} for an illustration):

\begin{enumerate}
	\item \emph{Path merging}: Two paths can be merged into one longer path if $H(V,D)$ has an arc from the end of one path to the start of the other path.
	
	\item \emph{Path unbalancing}: Consider any two paths $P$ and $P'$ with $1 < |P| \le |P'|$, let $v_s$ and $v_t$ denote the first and last vertices of $P$, and let $v'_s$ and $v'_t$ denote the first and last vertices of $P'$.  If $(v_s,v'_s)\in D$ we may remove $v_s$ from $P$ and append it at the beginning of $P'$. Likewise, if $(v'_t,v_t) \in D$ we may remove $v_t$ from $P$ and append it at the end of $P'$.
	
	\item \emph{Rotation}: if there is a path $P$ (say, $v_s, \ldots, v_t$) such that $(v_t, v_s)\in D$, we may add the arc $(v_t, v_s)$ to $P$ (obtaining a cycle), and then remove any other single arc from the resulting cycle to get a path $P'$. Observe that $P'$ and $P$ have the same vertex set, but they differ in their starting vertex along the cycle $v_s, \ldots, v_t, v_s$.
	
\end{enumerate}

\begin{figure}[h!]
	\begin{center}
\includegraphics[trim={0 0 0 0},clip,scale=.47]{figure3}
	\caption{(a) A spoiling graph $H$; (b) Path cover 1 --- a path cover of $H$; (c) Path cover 2, obtained from Path cover 1 by \emph{merging} paths $P_2,P_3$; (d) Path cover 3, obtained from Path cover 2 by \emph{unbalancing} paths $P_2,P_3$. One can verify that Path cover 3 is \emph{maximal}. In Path cover 3, $V_1=\{v_4\}$, $S=\{v_3,v_2\}$, $T=\{v_6,v_7\}$, $I=\{v_1,v_8,v_5\}$.
			To illustrate the notion of \emph{rotation}, suppose $H$ had contained two additional edges $(v_7,v_2), (v_1,v_3)$, then Path cover 3 would not have been maximal, as one could \emph{rotate} $P_3$ to get $P'_3 = v_8,v_5,v_7,v_2,v_1$, then merge it with $P_2$.}
		\label{fig:path-cover}
	\end{center}
\end{figure}

A path cover is \emph{maximal} if no path merging operation and no path unbalancing can be applied to it, and furthermore, this continues to hold even after performing any single rotation operation.
%

\begin{proposition}
	\label{pro:maximalcover}
	Given a bipartite graph $G(U,V;E)$ with a perfect matching $\{(u_i,v_i)\}$, a maximal path cover in the associated spoiling graph $H(V,D)$ can be found in $O(n^3)$ time.
\end{proposition}

\begin{proof}
	Start with the trivial path cover in which each vertex of $V$ forms a path of length~1, and perform arbitrary path merging and path unbalancing operations (some of which are preceded by a single rotation operation) until no longer possible. The process must end within $O(n^2)$ operations, because each path merging and each path unbalancing operation increases the sum of squares of the lengths of the paths, and the sum of squares of the lengths is at most $n^2$. Each of these operations can be performed in $O(n)$ time.
\end{proof}

Given a maximal path cover %
of $H$ (where $p$ denotes the number of paths in the path cover), sort the paths in order of increasing lengths, breaking ties arbitrarily. Hence $1 \le |P_1| \le |P_2| \le \ldots \le |P_p|$. We consider the following classes of vertices of $V$:

\begin{enumerate}
	
	\item \emph{Singleton vertices} $V_1$. These are the vertices that belong to paths of length~1 in the given maximal path cover. Let $k = |V_1|$ denote the number of singleton vertices. Observe that $|P_k| = 1$ and $|P_{k+1}| > 1$.
	
	\item \emph{Other vertices} $V_2 = V \setminus V_1$. We partition $V_2$ into three subclasses of vertices:
	
	\begin{enumerate}
		
		\item \emph{Start vertices} $S$. These are the starting vertices of those paths that have length larger than~1. The start vertex of path $j$, for $k < j \le p$, is denoted by $s_j$.
		
		\item \emph{End vertices} $T$. These are the end vertices of those paths that have length larger than~1. The end vertex of path $j$, for $k < j \le p$, is denoted by $t_j$.
		
		\item \emph{Intermediate vertices} $I = V_2 \setminus (S \cup T)$.
		
	\end{enumerate}
	
\end{enumerate}



\begin{lemma}
	\label{lem:properties}
The classes of vertices listed above have the following properties. %
	\begin{enumerate}
		\item There are no arcs in $H$ between vertices of $V_1$. Hence no vertex of $V_1$ can be a spoiler vertex for another vertex in $V_1$.
		\item There are no arcs in $H$ from vertices of $V_1$ to vertices in $S$. Hence no vertex of $V_1$ can be a spoiler vertex for a vertex in $S$.
		\item There are no arcs in $H$ from vertices of $T$ to vertices in $V_1$. Hence no vertex of $T$ can be a spoiler vertex for a vertex in $V_1$.
		\item For $i \not= j$, there are no arcs in $H$ from any vertex $t_i\in T$ to any vertex $s_j \in S$. Hence no vertex of $T$ can be a spoiler vertex for a vertex in $S$, unless they both belong to the same path in the given maximal path cover.
		\item $(s_i,s_j) \not\in D$ for any $s_i,s_j \in S$ with $i < j$.  Hence $s_i$ cannot be a spoiler vertex for $s_j$ if $i < j$.
		\item $(t_j,t_i) \not\in D$ for any $t_i,t_j \in S$ with $i < j$.  Hence $t_j$ cannot be a spoiler vertex for $t_i$ if $i < j$.
		\item If for some $s_j \in S$ and $t_j \in T$ (where $s_j$ and $t_j$ are start and end vertices of the same path $P_j$) 
%
we have  %
$(t_j,s_j)\in D$, then there are no arcs from $(T \setminus \{t_j\}) \cup V_1$ to any of the vertices of $P_j$, and likewise, no arcs from $s_{k+1}, \ldots, s_{j-1}$ to any of the vertices of $P_j$.
	\end{enumerate}
\end{lemma}

\begin{proof}
	All properties follow from the maximality of the path cover. Properties 1,2,3 and 4 hold because otherwise one could perform a path merging operation. Properties 5 and 6 hold because otherwise one could perform a path unbalancing operation. Property 7 holds because otherwise one could perform a rotation operation for path $P_j$, followed either by a path merging operation (if there is an arc from $(T \setminus \{t_j\}) \cup V_1$ to any of the vertices of $P_j$) or a path unbalancing operation (if there is an arc from $s_{k+1}, \ldots, s_{j-1}$ to any of the vertices of $P_j$).
\end{proof}

%
%

%
%
%
%

We now introduce additional notation.
Considering only the arcs in $D$ leading from $V_2$ to $V_1$, we let $M_{21}$ denote the maximum matching among these arcs. In our analysis we shall consider three parameters:

\begin{enumerate}
	
	\item $\epsilon_1$: its value is such that $k = \left(1/2 - \epsilon_1\right)n$ (recall that $k = |V_1|$ is the number of singleton paths in the maximal path cover). Observe that $\epsilon_1$ might be negative.
	
	\item $\epsilon_2$: its value is such that $p = k + \epsilon_2 n =  \left(1/2 - \epsilon_1 + \epsilon_2 \right)n$ (recall that $p$ is the total number of paths in the maximal path cover). Necessarily, $\epsilon_2 \ge 0$.
	
	\item $\epsilon_3$: its value is such that $|M_{21}| = \left(1/2 - \epsilon_3\right)n$. Necessarily, $\epsilon_3 \ge 0$.
	
\end{enumerate}

Given the above classes of vertices, we consider four possible candidate linear orders for $\pi$ (denoted $\pi_1, \pi_2, \pi_3, \pi_4$, see below for details). Given some linear order $\pi$, we shall use the notation $\rho(\pi)$ to denote the fraction of vertices guaranteed to be matched under $\pi$. This fraction will be expressed as a function of the parameters $\epsilon_1$, $\epsilon_2$ and $\epsilon_3$, and we will show that regardless of the value of these parameters, there must be some $\pi$ with $\rho(\pi) \ge 22/43$.

The following four lemmas present the four candidate linear orders for $\pi$ along with their corresponding guarantees.
Whenever unspecified, the order within a set of vertices can be arbitrary; \eg, for two sets of vertices $X,Y$, $\pi = X,Y$ means that the set $X$ precedes $Y$ and the order within $X$, as well as the order within $Y$, is arbitrary.

\begin{lemma}
	\label{lem:smallM21}
	For $G$ and $\pi_1 = V_2,V_1$, $$\rho(\pi_1) \ge \frac{1}{n}\left(|V_1|+\frac{|V_2|}{2} - \frac{|M_{21}|}{2}\right) = \frac{1}{2} - \frac{\epsilon_1}{2} + \frac{\epsilon_3}{2}.$$%
\end{lemma}

\begin{proof}
	Let $\sigma$ be an arbitrary linear order over $U$. Let $m$ denote the number of vertices in $V_2$ that are matched to vertices in $U_1$, where $U_1$ is the set of partners of $V_1$.
	Then $m \le |M_{21}|$. Of the $|V_2| - m$ vertices of $V_2$ not matched to vertices in $U_1$, at least half are matched (because for every unmatched vertex from this set, its partner must be matched to a different vertex from this set). In addition, all those vertices of $V_1$ whose partner is not matched to $V_2$ are matched, because of property~1 of \expref{Lemma}{lem:properties}. Hence the total number of vertices matched is at least 
	\begin{eqnarray}
m + \frac{|V_2| - m}{2} + |V_1| - m \ge |V_1|+\frac{|V_2|}{2} - \frac{|M_{21}|}{2},
	\end{eqnarray}
	as desired.
\end{proof}

\begin{lemma}
	\label{lem:largeM21}
	For $G$ and $\pi_2 = V_1,V_2$,
$$\rho(\pi_2) \ge \frac{2}{3} - \frac{1}{3}(\epsilon_1 + \epsilon_3).$$
\end{lemma}

\begin{proof}
	%
	Let $\sigma$ be an arbitrary linear order over $U$. All vertices in $V_1$ are matched because of property~1 of \expref{Lemma}{lem:properties}. As to the vertices in $V_2$, observe that $|N(V_2)| \ge |V_2| + |M_{21}|$, as the set $N(V_2)$ includes the $|V_2|$ partners of $V_2$, plus at least $|M_{21}|$ additional neighbors in $U_1$ (due to the matching $M_{21}$). Moreover, if $x$ vertices are removed from $V_2$, the number of remaining neighbors is at least $|V_2| + |M_{21}| - 2x$, because each vertex of $V_2$ contributed at most two neighbors to the lower bound that we provided on the number of neighbors.
	
	Let $x$ denote the number of vertices in $V_2$ matched under $(\pi_2,\sigma)$.
	Then the size of the matching is $|V_1| + x$, the number of unmatched vertices in $V_2$ is $|V_2|-x$, and they have at least $|V_2| + |M_{21}| - 2x$ neighbors which have to be matched. Since the number of matched vertices on each side is the same, %
	we have that $|V_1| + x \ge |V_2| + |M_{21}| - 2x$.
	
	
	We get that
	\begin{eqnarray}
		3x \  \geq \  |V_2|+|M_{21}|-|V_1|
		& = & n\left(\frac{1}{2}+\epsilon_1\right)+n\left(\frac{1}{2}-\epsilon_3\right)-n\left(\frac{1}{2}-\epsilon_1\right)\\
		& = & \left(\frac{1}{2} + 2\epsilon_1 - \epsilon_3\right)n.
	\end{eqnarray}
	Therefore, the size of the matching is at least
	\begin{eqnarray}
	|V_1| + x \geq \left( \frac{1}{2} - \epsilon_1\right)n + \left(\frac{1}{6} + \frac{2\epsilon_1}{3} - \frac{\epsilon_3}{3}\right)n = \left(\frac{2}{3} - \frac{1}{3}(\epsilon_1 + \epsilon_3)\right) n.    \qedhere
	\end{eqnarray}
	%
\end{proof}

\begin{lemma}
	\label{lem:sort1}
For $G$ and $\pi_3 = t_p, \ldots, t_{k+1}, V_1, s_{k+1}, \ldots, s_p, I$,
$$\rho(\pi_3) \ge \frac{2p - k}{n} = \frac{1}{2} - \epsilon_1 + 2\epsilon_2.$$
	%
\end{lemma}

\begin{proof}
	In $\pi_3$, we refer to the vertices of $T \cup V_1 \cup S$ as the \emph{prefix} of $\pi_3$, and to the vertices of $I$ as the \emph{suffix}. \expref{Lemma}{lem:properties} implies that within the prefix, the only arcs of $H$ that go from an earlier vertex to a later vertex are of the form $(t_j,s_j)$ (for a path $P_j$ that can undergo a rotation). We claim that regardless of $\sigma$, all the prefix will be matched. As the length of this prefix is $2p - k$, this will prove the lemma.
	
	It remains to prove the claim. Suppose first that in the above prefix there are no arcs of $H$ that go from an earlier vertex to a later vertex. Then earlier vertices
	in this prefix cannot be spoiling vertices for later vertices. Hence indeed,
	regardless of $\sigma$, all the prefix will be matched.
	
	Suppose now that in the prefix of $\pi_3$ there are arcs of $H$ that go from an earlier vertex to a later vertex. As noted above, such an arc would be of the form $(t_j,s_j)$. We need to show that even if $t_j$ acts as a spoiling vertex for $s_j$ under $\pi_3$ and $\sigma$, the vertex $s_j$ will still be matched. Consider the path $P_j$, and let us rename its vertices as $x_1, \ldots, x_{\ell}$ (where previously we used $s_j = x_1$ and $t_j = x_{\ell}$). We wish to show the $x_1$ would be matched even if $x_{\ell}$ is matched to the partner of $x_1$. The path $P_j$ implies that the partner of $x_2$ is a neighbor of $x_1$ in $G$. Hence $x_1$ will be matched if no vertex preceding $x_1 = s_j$ in $\pi_3$ is matched to the partner of $x_2$. By property~7 of \expref{Lemma}{lem:properties}, there is no arc in $H$ from any of the vertices $T \cup V_1 \cap \{s_{k+1}, \ldots, s_{j-1}\} \setminus \{t_j\}$ to $x_2$, and consequently none of them can be matched to the partner of $x_2$. As to $t_j = x_{\ell}$, it might be a neighbor of the partner of $x_2$ (in fact, it could be that $\ell = 2$), but we already assumed that $t_j$ is matched to the partner of $x_1$, and hence it is not matched to the partner of $x_2$. Hence no vertex preceding $x_1 = s_j$ in $\pi_3$ is matched to the partner of $x_2$, and hence $s_j$ will be matched.
\end{proof}

Let $V_e$ ($V_o$, respectively) denote those vertices of $S \cup I$ that are at even (odd, respectively) distance from the beginning of their respective path. Observe that $S \subset V_e$.

\begin{lemma}
	\label{lem:sort2}
	For $G$ and $\pi_4 = t_p, \ldots, t_{k+1}, V_1, V_o, V_e$,
	$$\rho(\pi_4) \ge \frac{5}{9} - \frac{p}{9n} = \frac{1}{2} + \frac{\epsilon_1}{9} - \frac{\epsilon_2}{9}.$$
	%
\end{lemma}

\begin{proof}
	Observe that $|V_e| \ge |V_o|$, because in every path (of length above~1) the vertices alternate in entering $V_e$ and $V_o$, and start with $V_e$. Observe also that every vertex $v \in V_e$ contributes two distinct neighbors to $N(V_e)$: the partner of $v$, and the partner of the vertex that follows $v$ on its path (note that the vertex that follows $v$ is not in $V_e$). Likewise, every vertex $v \in V_o$ contributes two distinct neighbors to $N(V_o)$.
	
	%
	%
	%
	%
	
	%
	
	%
	
	%
	
	Regardless of $\sigma$, all $p$ vertices of $T$ and $V_1$ are matched, as in \expref{Lemma}{lem:sort1}. For a given $\sigma$, let $n_o$ be the number of vertices matched in $V_o$ and let  $n_e$ be the number of vertices matched in $V_e$.
	Then, $|V_o| - n_o$, the number of unmatched vertices  in $V_o$, satisfies $2(|V_o| - n_o) \le p + n_o$, because the neighbors of the unmatched vertices in $V_o$ need to be matched to earlier vertices in $T \cup V_1 \cup V_o$.
	Likewise,  $|V_e| - n_e$, the number of unmatched vertices  in $V_e$, satisfies $2(|V_e| - n_e) \le p + n_o + n_e$.
	Adding two times the first inequality and three times the second we get that $4|V_o| + 6|V_e| - 4n_o - 6n_e \le 5p + 5n_o + 3n_e$. Using $|V_o| + |V_e| = n-p$ and $|V_e| \ge |V_o|$, we can replace $4|V_o| + 6|V_e|$ by $5(n-p)$, implying that $9(p + n_o + n_e) \ge 5n - p$, as desired.
\end{proof}

%

We can now prove \expref{Theorem}{thm:86}.

\begin{proof}
	%
	Observe that $\rho(G) \ge \max_{i \in [1,4]}[\rho(\pi_i)]$.
	\begin{eqnarray}
		& & \mbox{By \expref{Lemma}{lem:smallM21} we have:} \quad \quad
		\rho(\pi_1) \ge \frac{1}{2} - \frac{\epsilon_1}{2}  + \frac{\epsilon_3}{2}.\\
		& & \mbox{By \expref{Lemma}{lem:largeM21} we have:} \quad \quad
		\rho(\pi_2) \ge \frac{2}{3} - \frac{1}{3}(\epsilon_1 + \epsilon_3).\\
		& & \mbox{By \expref{Lemma}{lem:sort1} we have:} \quad \quad
		\rho(\pi_3) \ge \frac{1}{2} - \epsilon_1 + 2\epsilon_2.\\
		& & \mbox{By \expref{Lemma}{lem:sort2} we have:} \quad \quad
		\rho(\pi_4) \ge \frac{1}{2} + \frac{\epsilon_1}{9}  - \frac{\epsilon_2}{9}.
	\end{eqnarray}
	
	
	Taking a weighted average of the lower bounds provided by the four lemmas, with weights $$\frac{2}{43}, \frac{3}{43}, \frac{2}{43}, \frac{36}{43},$$ respectively, results in a weighted average of $22/43$. Hence regardless, of the values of $\epsilon_1$, $\epsilon_2$ and $\epsilon_3$, at least one of the lemmas gives $\rho(G) \ge 22/43$. For $$\epsilon_1 = \frac{19}{86}, \epsilon_2 = \frac{10}{86}, \epsilon_3 = \frac{21}{86},$$ none of the lemmas implies a bound better than $$\frac{1}{2} + \frac{1}{86} = \frac{22}{43}.$$
	
	The above analysis leads to a polynomial-time algorithm for finding $\pi$ that ensures $\rho(G) \ge 22/43.$ A maximal path cover of $H(V,D)$ can be found in polynomial time by \expref{Proposition}{pro:maximalcover}. Thereafter, the sets $V_1$, $S$, $T$, $V_e$ and $V_o$ can easily be determined, and likewise, the values of $\epsilon_1$ and $\epsilon_2$ can be easily computed. A maximum matching $M_{21}$ (from $V_2$ to $V_1$ in $H$) can be computed in polynomial time using any standard algorithm for maximum bipartite matching. Thereafter, $\epsilon_3$ can be easily computed. Given the values $\epsilon_1, \epsilon_2$ and $\epsilon_3$, one can determine which of $\pi_1, \pi_2, \pi_3$ or $\pi_4$ provides a higher lower bound on $\rho$, and use that linear order as $\pi$.
\end{proof}




\section{Regular graphs}
\label{sec:regular}

In this section we consider the case where $G(U,V;E)$ is a $d$-regular bipartite graph with $2n$ vertices. Given that such graphs have $d$
%
edge-disjoint  %
perfect matchings, one can hope to achieve higher values for $\rho$ for these graphs.

\subsection{Positive result}
\label{sec:regular-pos}

The following known %
%
fact %
(see for example~\cite{NaorW15}) establishes a lower bound on $\rho$, as a function of $d$.
A proof is provided for completeness.

\begin{proposition}
	\label{pro:smalld}
For every $d$-regular graph $G \in G_n$, 
%
the inequality  $\rho[G] \ge  d/(2d-1)$ holds.  %
\end{proposition}

\begin{proof}
	Since the greedy algorithm produces a maximal matching, it suffices to show that every maximal matching in a $d$-regular graph has size at least $\frac{d}{2d-1}n$.
	To see this, let $S \subset U$ and $T \subset V$ be the sets of unmatched nodes in an arbitrary maximal matching, and suppose $|S|=|T|=(1-\alpha) n$. The nodes in $S,T$ must form an independent set.
	Consider the size of the edge set connecting $S$ and $V \setminus T$.
	On the one hand, this size equals $(1-\alpha) n d$ (since all edges from $S$ go to $V \setminus T$); on the other hand, this size is at most $\alpha n (d-1)$ (since at least one edge from each node in $V \setminus T$ goes to $U \setminus S$).
	Thus, $(1-\alpha) n d \leq \alpha n (d-1)$, implying that $\alpha \geq d/(2d-1)$.
	Hence we have that $|M_G[\sigma,\pi]| \ge  \frac{d}{2d-1}n$, for every $\pi$.
\end{proof}

\begin{remark}
	For every $d$ there exists a $d$-regular graph with a perfect matching that admits a maximal matching of size $\frac{d}{2d-1}n$. Suppose that $n = 2d-1$, and
consider a $d$-regular graph where $|S|=|T|=d-1$ for some $S \subset U, T \subset V$, every node in $U \setminus S$ is connected to a single, different node in $V \setminus T$, and to all $d-1$ nodes in $T$, and every node in $V \setminus T$ is connected to a single, different node in $U \setminus S$, and to all $d-1$ nodes in $S$. The perfect matching between $U \setminus S$ and $V \setminus T$ is a maximal matching of size $\frac{d}{2d-1}n$.
\end{remark}

The lower bound of \expref{Proposition}{pro:smalld} approaches
%
%
$1/2$ from above as $d$ grows.  
The following theorem shows that there exists %
%
%
%
%
a   %
linear order $\pi$ that ensures that the fraction of matched vertices approaches $5/9$.
This is a direct corollary from \expref{Lemma}{lem:sort2} and a theorem in \cite{FeigeRS14}.

%
\begin{corollary}
	\label{cor:regular}
	For $d$-regulars bipartite graphs, 
$$\rho \ge \frac{5}{9} - O\left(\frac{1}{\sqrt{d}}\right).$$	
\end{corollary}

\begin{proof}
  Theorem~3 in~\cite{FeigeRS14} shows that every $n$-vertex $d$-regular graph has a path cover (referred to as a \emph{linear forest}) with  $p = O(n/\sqrt{d})$ paths. This means that there exists a maximal path cover with $O(n/\sqrt{d})$ paths, as the operations used to reach a maximal path cover (path merging, path unbalancing and rotation) can only decrease the number of paths.
%
%
The desired lower bound on $\rho(G)$ follows by \expref{Lemma}{lem:sort2}. %
%
%
%
\end{proof}

\begin{remark}
\begin{enumerate}
	
	\item For small $d$, the bound of $\rho \ge d/(2d - 1)$ which holds for every maximal matching is stronger than the bound in \expref{Corollary}{cor:regular}.
	
	\item The proof of \expref{Corollary}{cor:regular} extends to graphs that are nearly $d$-regular, by using Theorem~5 from~\cite{FeigeRS14}.
	
	\item For $d$-regular graphs, conjectures mentioned in~\cite{FeigeRS14} combined with our proof approach suggest that $\rho \ge 5/9 - O(1/d)$.
	
\end{enumerate}
\end{remark}

\subsection{Negative result}
\label{sec:regular-neg}

The following example shows that even in a regular graph with arbitrarily high degree, there may be no linear order $\pi$ that ensures to match more than a fraction $8/9$ of the vertices.

\begin{theorem}\label{thm:reg_89}
	For every integers $d, t \ge 1$, there is a regular bipartite graph $G_{d,t}$ of even degree $2d$ and $n= 3dt$ vertices on each side such that $\rho(G_{d,t}) \le 8/9$.
\end{theorem}

\begin{proof}
	Consider a regular bipartite graph $G(U,V;E)$ with even degree $2d$, and $3d$ vertices on each side. To define the edge set, let $U = U_1 \cup U_2 \cup U_3$ with each $U_i$ of cardinality $d$, and similarly $V = V_1 \cup V_2 \cup V_3$ with each $V_i$ of cardinality $d$. For every $i \not= j$, we have a complete bipartite graph between $U_i$ and $V_j$, and for every $i$, there are no edges between $U_i$ and $V_i$.
	
	Let $\pi$ be an arbitrary linear order over $V$, let $S$ be the first $2d$ vertices in $\pi$, and let $T$ be the last $d$ vertices. Let $i$ be such that $|V_i \cap T|$ is largest (breaking ties arbitrarily). Without loss of generality we may assume that $i = 3$, and then $|V_3 \cap T| \ge d/3$. Hall's condition implies that there is a perfect matching between $U_1 \cup U_2$ and $S$ (and more generally, between $U_1 \cup U_2$ and any $2d$ vertices from $V$). Hence one can choose a linear order $\sigma$ over $U$ whose first $2d$ vertices are $U_1 \cup U_2$ that will match the vertices of $S$ one by one. Thereafter, the vertices of $T \cap V_3$ will remain unmatched.
	
	To get the graph $G_{d,t}$ claimed in the theorem, take $t$ disjoint copies of $G(U,V;E)$ above.
\end{proof}






\section{Random linear order}
\label{sec:random}

In this section we consider scenarios in which the maximizing player is unaware of the graph structure.
In such scenarios, the best she can do is impose a random linear order over the vertices in $V$.

We first show that there exists a graph $G \in G_n$ for which a random linear order does not match significantly more than %
half of the vertices.%

\begin{proposition}
	There exists a bipartite graph $G(U,V;E) \in G_n$ such that a random linear order gets $\rho(G)=1/2+o(1)$ almost surely.
	\label{ex:half}
\end{proposition}

\begin{proof}
	Consider the graph $G(U,V;E)$, where $U=(U_1,U_2)$, $V=(V_1,V_2)$, and each of $U_1,U_2,V_1,V_2$ is of size $n/2$. The set of edges constitutes a perfect matching between $U_1$ and $V_1$, a perfect matching between $U_2$ and $V_2$, and a bi-clique between $U_1$ and $V_2$.
	Let $\pi$ be a random linear order. Rename the vertices in $V_1$ such that $v_{1j}$ is the $j$-th highest vertex from $V_1$ in the linear order $\pi$, and $u_{1j}$ is $v_{1j}$'s match in the perfect matching between $U_1$ and $V_1$.
	Consider an arrival order $\sigma$ in which vertex $u_{1j}$ is the $j$-th vertex (and $U_2$ arrives after $U_1$ in an arbitrary order). All vertices of $U_1$ will be matched (as they arrive first, there are only $n/2$ of them, and each of them has more than $n/2$ neighbors in $V_2$). Vertices in $V_1$ can only be matched with vertices of $U_1$. 
	
	Let $k$ denote the number of vertices in $V_1$ that are matched, and let $v_{1\ell}$ be the highest index (\ie, lowest ranked) vertex in $V_1$ to be matched (that is, all vertices $v_{1j}$ with $j>\ell$ are unmatched). Necessarily, $v_{1\ell}$ is matched with $u_{1\ell}$. This can happen only if every vertex from $V_2$ that precedes $v_{1\ell}$ according to $\pi$ is matched at the time that $u_{1\ell}$ arrives. As $k-1$ vertices of $V_1$ are matched before $u_{1\ell}$ arrives, there can be at most $\ell - k$ vertices from $V_2$ that precede $v_{1\ell}$ in the linear order $\pi$. Hence, for $k$ vertices from $V_1$ to be matched, the linear order $\pi$ should be such that for some $\ell$ in the range $k \le \ell \le n/2$, in the prefix of $\pi$ of length $2\ell-k$, there are at least $\ell$ vertices from $V_1$ and at most $\ell - k$ vertices from $V_2$. Such an $\ell$ is said to be \emph{good}. 
	
	We claim that the probability that a particular $\ell$ is good is smaller than the probability that in $2\ell-k$ random tosses of an unbiased coin, at most $\ell-k$ tosses come up heads. To see this, note that the prefix of $\pi$ is sampled from $V$ without repetition; the random coin version is equivalent to sampling from $V$ with repetition, where in every sample, regardless of past samples, there is probability %
 $1/2$   %
of sampling a vertex from $V_1$ (coin coming up tails), and probability
%
$1/2$   %
of sampling a vertex from $V_2$ (coin coming up heads). In both sampling methods, the expectation of the number of vertices from $V_2$ in the prefix is 
%
$\ell - (k/2)$.   %
 The probability of deviating by more than 
%
$k/2$  %
from the expectation is larger when sampling with repetitions than when sampling without repetition. 
	
Now, by standard bounds on binomial coefficients (or alternatively, by standard
%
Bernstein bounds\footnote{These concentration inequalities, 
commonly misattributed to Chernoff, were published by 
Sergei N. Bernstein in 1924 and 1937. (See the Wikipedia
entry
\href{https://en.wikipedia.org/wiki/Bernstein_inequalities_(probability_theory)}{%
``Bernstein inequalities (probability theory).''})}
this probability is at most 
	\begin{eqnarray}
		\eee^{-k^2/2(2\ell-k)} \le \eee^{-k^2/4\ell}.
	\end{eqnarray}
	A union bound over all values of $\ell$ (in the range
%
$k \le \ell \le n/2$)   %
shows that if $k \ge 2\sqrt{n\log n}$, it is highly unlikely that a good $\ell$ exists. Consequently, with overwhelming probability, the size of the matching is $(1/2 +o(1))n$, whereas $OPT=n$.
\end{proof}



In the above example, if the vertices of $V$ with degree 1 are placed in the prefix of $\pi$, then the %
matching obtained  %
is optimal. This might suggest that prioritizing
%
low-degree   %
vertices in $\pi$ (and randomizing within sets of vertices of comparable degrees) leads to good performance.
However, the example above can be transformed into one where all vertices in $V$ have the same degree, and moreover, each vertex has a degree of $\Omega(\sqrt{n})$.
To see this, consider a graph where vertices are partitioned into sets of perfect matchings of size $\sqrt{n}$, $\{(U_{11},V_{11}),\ldots, (U_{1\sqrt{n}},V_{1\sqrt{n}}), (U_{21},V_{21}),\ldots, (U_{2\sqrt{n}},V_{2\sqrt{n}})\}$. Each $V_{1i}$ is also connected in a bi-clique to $U_{2i}$, and in addition, there are sets $U',V'$ of size $\sqrt{n}$ each connected to the vertices of the other side to balance out the degrees. A similar argument shows that in this graph, a random linear order performs badly as well. 

In contrast to the last examples, in some classes of graphs, a random linear order guarantees to match a fraction of the vertices that is bounded away from a half.
This is the case, for example, in Hamiltonian graphs.
The formal statement and proof are deferred to \expref{Section}{sec:hamiltonian}.


%
\section{Finding a perfect \texorpdfstring{$\pi$}{pi}}   %
%
\label{sec:perfect-pi}

%

A linear order $\pi$ over $V$ is said to be \emph{perfect} if for every linear order $\sigma$ over $U$, $|M_G[\sigma,\pi]| = n$.
A vertex $v \in V$ is \emph{good} with respect to a set of vertices $S \subset V$ if there is no perfect matching between $N(v)$ and $S$.
Given a linear order $\pi$ over $V$, let $\precvpi$ be the set of vertices preceding $v$ in $\pi$.

\begin{observation}
	\label{obs:perfect-pi}
	$\pi$ is perfect if and only if every vertex $v \in V$ is good with respect to $\precvpi$.
\end{observation}
\begin{proof}
	Suppose there exists a vertex $v \in V$ that is not good. Then, consider a linear order $\sigma$ where the vertices in $N(v)$ are placed first, in an order corresponding to the rank (according to $\pi$) of their partners in the perfect matching between $N(v)$ and $\precvpi$. In such a $\sigma$, $v$ will not be matched.
	Now suppose that all vertices in $V$ are good.
	Then, for every $\sigma$, for every $v \in V$, there exists a vertex $u \in U$ that is not matched to a vertex in $\precvpi$; therefore $v$ will surely be matched.	
\end{proof}


We present an algorithm that finds a perfect $\pi$ if one exists, and claims that no such $\pi$ exists otherwise.

Checking whether $v_i$ is good can be done in polynomial time by running a maximum matching algorithm on $N(v_i)$ and $S_i$.

\medskip
\noindent\textbf{An algorithm for finding a perfect $\pi$:}
\begin{enumerate}
	\item Let $S_1=V$.
%
%
\item For $i=1,\ldots,n$:    %
	\begin{enumerate}
		\item Find an arbitrary good vertex $v_i \in S_i$ with respect to $S_i \setminus \{v_i\}$, and place it in rank $n-i+1$ in $\pi$.
		\item Set $S_{i+1}=S_i \setminus \{v_i\}$.
	\end{enumerate}
\end{enumerate}

\begin{lemma}
	If there exists a perfect $\pi$, then the algorithm is guaranteed to find a good $v_i$ in every iteration $i$.
\end{lemma}
\begin{proof}
	Consider some perfect linear order $\pi$ (not necessarily the one produced by our algorithm), and the suffix $V_{i-1}=v_{i-1},v_{i-2},\ldots,v_1$ of vertices chosen in the first $i-1$ iterations of the
algorithm. %
(Of course, there must be a good $v_1$ at the first iteration, otherwise there is no perfect $\pi$.)  %
 Let $\pi'$ be the linear order that places $v_{i-1},v_{i-2},\ldots,v_1$ as the lowest-ranked %
vertices in the same order as the algorithm picked them, and places all other vertices of $V\setminus V_{i-1}$ in a higher rank than $V_{i-1}$ according to their internal order in $\pi$.
	
Since every $v_j$, $1\leq j\leq i-1$, is good with respect to $V\setminus \{v_1,\ldots, v_{j}\}$, %
clearly $v_j$ is good with respect to $\overline{\pi'}(v_j)$ (since $\overline{\pi'}(v_j) = V\setminus \{v_1,\ldots, v_{j}\}$). Now consider a vertex $v\in V\setminus V_{i-1}$. This vertex is good with respect to $\precvpi$, and since $\overline{\pi'}(v)\subseteq  \precvpi$, it is clear that $v$ is good with respect to $\overline{\pi'}(v)$. It follows that $\pi'$ is perfect as well.
	
	Let $v'_{i}$ be the vertex ranked in position $n-i+1$ in $\pi'$. Since $\pi'$ is perfect, this vertex is good with respect to $\overline{\pi'}(v'_i)$. But since $\overline{\pi'}(v'_i)\cup \{v'_{i}\}$ is exactly the set $S_i$ in iteration $i$, it is guaranteed that the algorithm can find a good $v_i$ in this iteration.
\end{proof}

Let $\pi$ be the linear order computed by the algorithm. Since every vertex $v$ is good with respect to $\precvpi$, it follows from \expref{Observation}{obs:perfect-pi} that $\pi$ is perfect, and \expref{Proposition}{thm:perfect-pi} follows. 


\section{Hamiltonian bipartite graphs}
\label{sec:hamiltonian}

In this section we establish two results about Hamiltonian graphs.
First, we show that %
$\rho \ge 5/9$.   %
Note that, since for the case of a Hamiltonian graph there exists a path cover using only a single path
%
(\ie, $p=1$), \expref{Lemma}{lem:sort2} directly implies that 
%
$\rho \ge 5/9 - 1/9n$.
\expref{Theorem}{thm:hamiltonian-5-9} improves this bound to $5/9$.
Second, we show that for Hamiltonian graphs, even a random linear order $\pi$ ensures a ratio that is bounded away from %
$1/2$.  %
(This is in contrast to general graphs, see \expref{Section}{sec:random}.) %

\begin{theorem}
%
%
%
%
  For Hamiltonian graphs, $\rho \ge 5/9$.
%
	\label{thm:hamiltonian-5-9}
\end{theorem}

\begin{proof}
Consider a Hamiltonian graph $G$ and a Hamiltonian cycle $u_1,v_1,u_2,v_2,\ldots, u_n,v_n,u_1$ 
%
%
in $G$ that passes through the vertices in the order listed. %
Let $V_o=\{v_i\ :\ i=2\ell+1,\ \ell\in \mathbb{N},\ i\leq n \}$ be the set of odd labeled vertices of the cycle, and $V_e=V\setminus V_o$.
	
	We first claim that if the number of vertices is even 
%
($|V_o|=|V_e|=n/2$),   %
then $\pi=V_e, V_o$ (and in fact, also $\pi=V_o, V_e$) ensures that $\rho\ge 5/9$.
%
Let $n_e$ and $n_o$ be the number of vertices of $V_e$ and $V_o$, respectively,
%
matched in $M_G[\sigma,\pi]$ defined using $\pi=(V_e,V_o)$ and an arbitrary $\sigma$. %
Similarly   %
to the proof of \expref{Lemma}{lem:sort2}, it is easy to see that each vertex in $V_e$ contributes two distinct neighbors to $N(V_e)$, and each vertex in $V_o$ contributes two distinct neighbors to $N(V_o)$. %
%
(The difference from the proof of \expref{Lemma}{lem:sort2} is that this property also holds for $v_1\in V_o$ and $v_n\in V_e$, and this follows because $v_1$ and $v_n$ contribute $u_1$ to $N(V_o)$ and $N(V_e)$, respectively.)  %
 The number of unmatched vertices in $V_e$, namely $|V_e|-n_e$, satisfies \begin{eqnarray}2(|V_e|-n_e)\leq n_e,   %
	\end{eqnarray} because the neighbors of the unmatched vertices in $V_e$ must be matched to vertices in $V_e$, as they precede the vertices in $V_o$ in  $\pi$. Likewise, the number of unmatched vertices in $N_o$, namely $|V_o|-n_o$, satisfies \begin{eqnarray}2(|V_o|-n_o)\leq n_e+n_o.
\end{eqnarray} Adding two times the first inequality and three times the second, we get \begin{eqnarray}4|V_e| + 6|V_o|\leq 9(n_o+n_e)\ \Rightarrow \ \frac{5}{9}\cdot n\leq n_o+n_e.
\end{eqnarray} As $|V|=n$ and  $|M_G[\sigma,\pi]|=n_o+n_e$, this implies that 
%
$\rho\geq 5/9.$   %
	
	We now handle the case where $n$ is odd. \expref{Lemma}{lem:sort2} ensures that $\frac{5}{9}\cdot n -\frac{1}{9}$ of the vertices are matched by $\pi_4$ when the path cover is of a single path.
	If $\frac{5}{9}\cdot n -\frac{1}{9}$ is not integral, then $\lceil \frac{5}{9}\cdot n -\frac{1}{9}\rceil$ is at least $\frac{5}{9}\cdot n$, thus $\rho \geq \frac{5}{9}$.
	Therefore, it only remains to handle the case where $\frac{5}{9}\cdot n -\frac{1}{9}$ is integral;
	%
	namely where $n=18\ell+11$ for some integer $\ell$ (recall that $n$ is odd in this case).
	In this case, we show that $\pi=V_e, V_o$ ensures that $|M[\sigma,\pi]|> \frac{5}{9}\cdot n$ for every $\sigma$. Since $n=18\ell+11$,
%
  it follows that $|V_o|=9\ell+6$ and $|V_e|=9\ell+5$. %
%
%
	As in the case where $n$ is even, every vertex in $V_e$ contributes two distinct neighbors to $N(V_e)$.
	As for $V_o$, every vertex in $V_o \setminus \{v_1,v_n\}$ also contributes two distinct neighbors to $N(V_o)$, and $v_1$ and $v_n$ contribute (together) to $N(V_o)$ three additional distinct vertices (since they share a vertex along the Hamiltonian cycles).
	Using the same reasoning as before, it follows that 
	\begin{eqnarray}
		2(|V_e|-n_e)\leq n_e\Rightarrow 18\ell+10\leq 3\cdot n_e\Rightarrow n_e\geq 6\ell+3\frac{1}{3}.
		\end{eqnarray}
	 Since $n_e$ is integral, this implies that
	\begin{eqnarray}
	n_e\geq 6\ell+4.
	\label{eq:eq1-hamilton}
	\end{eqnarray}
	Again, for $|V_o|-n_o$ we have $2(|V_o|-n_o)-1\leq n_o+n_e$. Rearranging gives us 
	\begin{eqnarray}
		n_e+3n_o\geq 18\ell+11.\end{eqnarray}
	Adding twice \expeqref{Inequality}{eq:eq1-hamilton} to the last inequality yields \begin{eqnarray}
		n_e+n_o\geq 10\ell+6\frac{1}{3}>10\ell +6\frac{1}{9}=\frac{5}{9}\cdot |V|,\end{eqnarray} which implies $\rho >\frac{5}{9}$. %
	This concludes the proof.
\end{proof}

Next, we show that for the case of a Hamiltonian graph, a random linear order yields a $\rho>1/2$.

\begin{theorem}
  Consider choosing a linear order $\pi$
%
%
%
%
%
%
%
%
%
%
at random. For every Hamiltonian graph $G$,
%
we have  %
$E_{\pi}[\min_{\sigma}[|M_G[\sigma,\pi]|]] > 0.5012$.
	\label{thm:random-hamiltonian}
\end{theorem}	

%
Since the proof is quite involved, we first give a proof overview before providing a formal proof. 

%
%
%
%

\subsection{Proof overview}
\label{sec:proof-overview}

We first provide a high-level %
overview of our proof.  %
%

A linear order $\pi$ (over $V$) is said to be \emph{safe} for a set $S \subset V$ if for every linear order $\sigma$ (over $U$) the greedy process matches at least one vertex in $S$ (\ie, no $\sigma$ leaves all vertices in $S$ unmatched).
Fix some constant $\epsilon$. In order to establish that $\rho \geq (1/2 + \epsilon)$, we need to show that there exists a linear order $\pi$ that is safe for every set $S$ of size $(1/2-\epsilon)n$.
Our proof approach is the following: we show that for a linear order $\pi$ chosen
%
at random, the expected number (expectation taken over choice of $\pi$) of sets of size $(1/2-\epsilon)n$ for which $\pi$ is unsafe is smaller than $1$.
This implies that there exists a linear order $\pi$ that is safe for all sets of size $(1/2-\epsilon)n$, as desired.

First, we define a collection of sets that can potentially remain unmatched (``bad'' sets).
%
%
%
Let $B_\epsilon$ denote the set of all sets $S \subset U$ of size $(1/2-\epsilon)n$ such that there exists a linear order $\pi$ that is unsafe for $S$.

Second, for a given set $S$ and a linear order $\pi$ we identify a sufficient condition for $\pi$ to be safe for $S$.
Let $S' \subset S$ be the lowest $\alpha n$ vertices in $S$ (according to $\pi$), let $v'$ be the last vertex in $S'$ (\ie, the vertex with rank $\alpha n$ in $S'$), and let $P$ be the set of vertices in $V - S'$ that precede $v'$ in $\pi$.
%
We claim that if the size of $P$ is smaller than the size of $N(S')$ (the neighbors of $S'$), then $\pi$ is safe for $S$.
To see this, assume by way of contradiction that $\pi$ is unsafe for $S'$.
This implies that every vertex in $N(S')$ is matched to a vertex in $V - S'$.
Since there are strictly less than $|N(S')|$ vertices in $V - S'$ that precede $v'$, at least one of the vertices in $N(S')$ must be matched to a vertex higher than $v'$. But, this vertex has a neighbor in $S'$ with lower rank, contradicting the greedy process.

We now proceed by establishing the following three lemmas:
\begin{itemize}
	\item \emph{Few bad sets lemma:} the size of $B_\epsilon$ is at most $n_B=n_B(\epsilon)$.
	\item \emph{Expansion lemma:}
	given a set $S \subset V$ and parameters $\alpha, \beta$, the probability (over a random choice of $\pi$) that the lowest $\alpha n$ vertices in $S$ have less than $\beta n$ neighbors is at most $p=p(\alpha,\beta)$.
	\item \emph{Good order lemma:} given a set $S \subset V$ and parameters $\alpha, \beta$, the probability (over a random choice of $\pi$) that the $(\alpha n)^{th}$ lowest vertex in $S$ is higher than $\beta n$ vertices in $V \setminus S$ is at most $q=q(\alpha,\beta)$.
\end{itemize}

The three lemmas are combined as follows.
For a given set $S$, due to the sufficient condition identified above, it follows from the union bound that the probability that a
%
random linear order $\pi$ is unsafe for $S$ is at most $p+q$.
Applying the union bound once more over all bad sets (at most $n_B$ sets, as implied by the few bad sets lemma), implies that the probability that a
%
random linear order $\pi$ is unsafe for some set of size $(1/2-\epsilon)n$ is at most $n_B(p+q)$.
Thus, to conclude the proof, it remains to find parameters such that $n_B(p+q) < 1$.

The good order lemma is independent of the graph structure.
In contrast, the expansion lemma and the few bad sets lemma rely heavily on the structure of the graph.
As it turns out, Hamiltonian graphs have properties that enable us to establish the two lemmas with good parameters.


\subsection{Full proof of \expref{Theorem}{thm:random-hamiltonian}}\label{sec:app-hamiltonian}


We use $H(\cdot)$ to denote the binary entropy function, %
\ie, given $p\in (0,1)$, \begin{eqnarray}H(p)=-p\log_2 p-(1-p)\log_2(1-p).
\end{eqnarray}

\begin{fact}[Stirling's Approximation]
	As $n\rightarrow\infty$, 
 %
 %
\begin{eqnarray}n!=(1+o(1))\sqrt{2\pi n}\left(\frac{n}{\eee}\right)^n.
		\end{eqnarray}
\end{fact}
Using Stirling's Approximation, one can derive the following bound.
\begin{fact}\label{fact:binom_estimate}
  For $n$ and $k=pn$ for some
constant %
%
$p\in (0,1)$, as $n\rightarrow\infty$,
	\begin{eqnarray}
	\binom{n}{k}=2^{(H(p)+o(1))n}.\label{eq:binom_approx}
	\end{eqnarray}
	%
\end{fact}

We first establish the good order lemma, which is independent of the graph structure.

\begin{lemma} [Good order lemma]
	Let $\alpha < \beta < 1$, $\rho= \frac{1}{2}+\epsilon$ for some $\epsilon>0$ %
	and $\bar{\rho}=1-\rho$ such that $\frac{\beta}{\alpha} > \frac{\rho}{\bar{\rho}}$. Let $S\subset V$ be a set of size $\bar{\rho}n$. The probability that in a random linear order $\pi$ there are at least $\beta n$ vertices of $V\setminus S$  before $\alpha n$ vertices from $S$ is at most $$2^{-(H(\alpha+\beta)-H(\frac{\alpha}{\bar{\rho}})\bar{\rho}-H(\frac{\beta}{{\rho}}){\rho}-o(1))n}.$$\label{lem:order}
\end{lemma}
\begin{proof}
	We first analyze the case that in the first $(\alpha+\beta)n$ vertices in $\pi$ there are \emph{exactly} $\alpha n$ vertices from $S$. The number of possibilities for this case is $\binom{\bar{\rho} n}{\alpha n}\binom{\rho n}{\beta n}$.
	
	Let $\beta' = \beta+x$ and $\alpha'=\alpha-x$. By the conditions on $\alpha,\beta$ and $\epsilon$, we have that $\frac{\beta'}{\alpha'}\geq\frac{\beta}{\alpha}\geq \frac{\rho}{\bar{\rho}}$. Therefore,
	\begin{eqnarray}
		& &  \beta'\bar{\rho}\geq \alpha'\rho\Rightarrow \beta'\bar{\rho}-\alpha'\beta'\geq \alpha'\rho-\alpha'\beta'\Rightarrow \frac{(\bar{\rho}-\alpha')}{\alpha'}\cdot \frac{\beta'}{(\rho-\beta')}\geq 1 \\
		& & \Rightarrow \frac{(\bar{\rho}n-\alpha'n+1)}{\alpha'n}\cdot \frac{\beta'n+1}{(\rho n-\beta'n)} > 1 \iff\frac{\binom{\bar{\rho} n}{\alpha' n}\binom{\rho n}{\beta' n}}{\binom{\bar{\rho} n}{\alpha' n-1}\binom{\rho n}{\beta' n+1}}>1.
	\end{eqnarray}
	It follows that
	%
	%
	$\binom{\bar{\rho}n}{\alpha n}\binom{\rho n}{\beta n}> \binom{\bar{\rho} n}{\alpha'n}\binom{\rho n}{\beta' n}$ for every $\alpha'< \alpha$  and $\beta'>\beta$ such that $\alpha+\beta=\alpha'+\beta'$. Therefore, the probability to have at most $\alpha n$ vertices from $S$ in the first $(\alpha+\beta)n$ vertices in $\pi$ is at most
	\begin{eqnarray}
		\frac{\alpha n\cdot \binom{\bar{\rho}n}{\alpha n}\binom{\rho n}{\beta n}}{\binom{n}{(\alpha+\beta)n}}&=&\frac{2^{(H(\frac{\alpha}{\bar{\rho}})\bar{\rho}+H(\frac{\beta}{{\rho}}){\rho}+o(1))n}}{2^{(H(\alpha+\beta)+o(1))n}}\\
		& = & 2^{-(H(\alpha+\beta)-H(\frac{\alpha}{\bar{\rho}})\bar{\rho}-H(\frac{\beta}{{\rho}}){\rho}-o(1))n},
	\end{eqnarray}
where the first equality follows
by  %
\expref{Fact}{fact:binom_estimate}.
\end{proof}

Let %
$\rho= 1/2 +\epsilon$  %
for some constant $\epsilon>0$,
and %
%
$\bar{\rho}=1-\rho= 1/2-\epsilon$. 
The next lemma will be used in order to prove the few bad sets lemma and the expansion lemma. It uses the existence 
of  %
a Hamiltonian cycle in the graph in order to claim that most sets will have a large number of neighbors. Therefore, a random set will have a large expansion. In addition, there will be few sets of size 
%
$(1/2-\epsilon)n$    %
with %
fewer than %
%
$(1/2+\epsilon)n$   %
neighbors (\ie, %
%
few bad sets).

\begin{lemma}
	Let $\alpha\in (0, 1/2)$ and $\beta \in (\alpha, 1)$ be two constants such that $\delta=\beta-\alpha < \alpha/2$. The number of sets $S$ of size $\alpha n$ where $|N(S)|\leq \beta n$ is at most $$2^{(\alpha H(\delta/ \alpha)+(1-\alpha)H(\delta/(1-\alpha))+o(1))n}.$$
	\label{lem:ham_alpha_beta}
\end{lemma}
\begin{proof}
	%
  Consider a Hamiltonian cycle $H=(v_1,u_1,v_2,u_2,\ldots, v_n,u_n,v_1)$, where $\{v_i\}_{i\in [n]}=V$ and $\{u_i\}_{i\in [n]}=U$. Let $S$ be
%
a  %
set of vertices from $V$ of cardinality $\alpha n$. Note that in the cycle $H$, each vertex $v$ of $S$ has two neighbors, where one of these neighbors is joined with  an adjacent vertex from $V$ in the cycle. Therefore, the number of neighbors of a sequence of $k$ consecutive vertices of $V$ in $H$ is $k+1$. Thus, the set $N(S)$ is of size $\alpha n$ plus the number of consecutive blocks of vertices from $V$ chosen.
	%
	
	We bound the number of ways to pick at most $\delta n$ consecutive blocks of vertices from $V$. We first bound the number of ways to pick \emph{exactly} $\delta n$ such blocks.
	In this case, the $\alpha n$ chosen elements have to be within $\delta n$ blocks. The number of ways to partition $\alpha n$ elements to $\delta n$
non-empty %
blocks is $\binom{\alpha n - 1}{\delta n - 1}$. After deciding the number of elements in each block, we need to figure out their  location along the Hamiltonian cycle. $(1-\alpha)n$ elements reside outside of the blocks of 
the  %
chosen $\alpha n$ elements. We need to %
choose  %
the location of the first block in $H$ (for which there are $n$ possibilities), and then the number of %
elements   %
%
between each pair of consecutive blocks %
where two blocks are separated by at least one element. The latter is equivalent to %
%
dividing %
$(1-\alpha)n$ elements into $\delta n$ 
non-empty  %
%
%
segments   %
for which there are $\binom{(1-\alpha)n-1}{\delta n - 1}$ possibilities. Overall, there are $n\binom{\alpha n - 1}{\delta n - 1}\binom{(1-\alpha)n-1}{\delta n - 1}$ such possibilities\footnote{Notice
%
%
there is  %
some %
overcounting  %
in this argument, but this bound suffices for our purpose.}.
	
	For $\delta'< \delta$, one can similarly prove the bound of 
	\begin{eqnarray}
		n\binom{\alpha n - 1}{\delta' n - 1}\binom{(1-\alpha)n-1}{\delta' n - 1}
	\end{eqnarray}
	which is smaller than 
	\begin{eqnarray}
		n\binom{\alpha n - 1}{\delta n - 1}\binom{(1-\alpha)n-1}{\delta n - 1}
	\end{eqnarray} by our conditions on $\alpha$ and $\delta$. Overall, we can bound the number of ways to pick at most $\delta n$ consecutive blocks of vertices from $V$ by
	\begin{eqnarray}
		\delta n^2\binom{\alpha n - 1}{\delta n - 1}\binom{(1-\alpha)n-1}{\delta n - 1} & < & \delta n^2\binom{\alpha n}{\delta n}\binom{(1-\alpha)n}{\delta n}\\
		&=& 2^{o(1)n}\cdot 2^{(H(\frac{\delta}{ \alpha})+o(1))\alpha n}\cdot 2^{(H(\frac{\delta} {(1-\alpha)})+o(1))(1-\alpha) n}\\
		& = &2^{(\alpha H(\frac{\delta}{ \alpha})+(1-\alpha)H(\frac{\delta} {(1-\alpha)})+o(1))n},
	\end{eqnarray}
	where the first equality follows \expref{Fact}{fact:binom_estimate}.
	%
\end{proof}

The expansion and few bad sets lemmas are obtained as direct corollaries of \expref{Lemma}{lem:ham_alpha_beta}.

%
\begin{lemma}[Few bad sets lemma for Hamiltonian graphs]
	Let $\epsilon$ be a constant such that $\epsilon<0.1$. The number of bad sets in any Hamiltonian graph is at most
	$$|B_\epsilon|\leq  2^{\left(\bar{\rho} H(\frac{2\epsilon}{\bar{\rho}})+\rho H(\frac{2\epsilon} {\rho})+o(1)\right)n}.$$ \label{lem:ham_bad_sets}
\end{lemma}
\begin{proof}
	Notice that if a set $S$ of size 
%
$\bar{\rho} n = (1/2-\epsilon)n$   %
has more than $\rho n$ neighbors, it cannot be left unmatched, since at least one of its neighbors will not be matched to $V\setminus S$. 
%
%
%
%
By \expref{Lemma}{lem:ham_alpha_beta}, the number of such sets is at most $2^{\left(\bar{\rho} H(\frac{2\epsilon}{\bar{\rho}})+\rho H(\frac{2\epsilon} {\rho})+o(1)\right)n}.$
\end{proof}

We note that this lemma is not true for general graphs.
An example of a graph that admits $2^{n/4}$ bad sets is given in \expref{Proposition}{ex:many-bad}.

\begin{lemma}[Expansion Lemma for Hamiltonian graphs]
	Consider a set $S\subset V$ of size $\bar{\rho} n$ and parameters $\alpha$, $\beta$. The probability that the lowest $\alpha n$ vertices in $S$ have
%
fewer  %
than $\beta n$ neighbors is at most
	$$2^{\left(-H(\frac{\alpha}{\bar{\rho}})+\alpha H(\frac{\delta}{ \alpha})+(1-\alpha)H(\frac{\delta} {(1-\alpha)})+o(1)\right)\bar{\rho}n}.$$   \label{lem:ham_expansion}	
\end{lemma}
\begin{proof}
	Consider a set $S$ of size $\bar{\rho} n$, and the first $\alpha n$ vertices in $S$ in a random linear order. This set is just a random subset of $S$ of size $\alpha n$. The number of choices of such subset is $$\binom{\bar{\rho}n}{\alpha n} = 2^{(H(\frac{\alpha}{\bar{\rho}})+o(1))\bar{\rho} n}.$$

 %
Notice that we can apply \expref{Lemma}{lem:ham_alpha_beta} 
 with set $S$ and bound the number of subsets of $S$ with ``small'' expansion, 
%
%
%
%
even though $S$ is a subset of $V$, because the same proof applies
%
%
with respect to a subset of vertices in only %
one side of a Hamiltonian graph. Therefore, the number of subsets of size $\alpha n$ of $S$ with at most $\beta n$ neighbors is at most
	\begin{eqnarray*}
		2^{(\alpha H(\frac{\delta}{ \alpha})+(1-\alpha)H(\frac{\delta} {(1-\alpha)})+o(1))\bar{\rho}n}.
	\end{eqnarray*}
	Combining the above, we get that the probability that a random set of $\alpha n$ vertices of $S$ have at most $\beta n$ neighbors is at most
	\begin{eqnarray*}
2^{(-H(\frac{\alpha}{\bar{\rho}})+\alpha H(\frac{\delta}{ \alpha})+(1-\alpha)H(\frac{\delta} {(1-\alpha)})+o(1))\bar{\rho}n}.     \qedhere
	\end{eqnarray*}          
\end{proof}

Now that we have established the three lemmas, %
we are ready to prove \expref{Theorem}{thm:random-hamiltonian}.

\begin{proof}[Proof of \expref{Theorem}{thm:random-hamiltonian}]
	Setting $\epsilon=0.0012$, $\alpha=0.245$ and $\beta=0.3675$ (and $\rho=\frac{1}{2}+\epsilon$, $\bar{\rho}=1-\rho$), we get that these parameters satisfy the conditions for \expref{Lemmas}{lem:order}, \ref{lem:ham_expansion} \expref{and}{lem:ham_bad_sets}.
	
	Applying \expref{Lemma}{lem:ham_bad_sets}, we get that the size of $B_\epsilon$ is at most $n_B\leq 2^{0.044n}$. Applying \expref{Lemma}{lem:ham_expansion}, we get that the probability that the lowest $\alpha n$ vertices of a set of size $\bar{\rho}n$ have %
fewer  %
than $\beta n$ neighbors is at most $p\leq 2^{-0.86n}$. Applying \expref{Lemma}{lem:order}, we get that the probability that for a set $S$ of size $\bar{\rho}n$ the $\alpha n$-th vertex in a random $\pi$ comes after $\beta n$ vertices of $V-S$ is at most $q\leq 2^{-0.45n}$. Combining these three
bounds, %
we get that the probability there exists a set of size $\bar{\rho}n$ unmatched by a random $\pi$ is at most $n_B(p+q)<1$, therefore, there must be a $\pi$ that matches at least one vertex in each set of size  $\bar{\rho}n$, and the proof follows.
\end{proof}	


%
%
This approach can also be   %
used to show that a random linear order
%
is guaranteed to  %
to match more than %
half of the vertices in every regular graph.
On the other hand, \expref{Theorem}{thm:3-4} in \expref{Section}{sec:additional} shows that one cannot hope to get $\rho > 3/4$ with a random linear order in regular graphs.

\begin{figure}[h!]
	\begin{center}
\includegraphics[trim={0 4.5cm 0 4cm},clip,scale=.5]{figure4}
\caption{The graph $G'$ admits two bad sets, namely $\{v_1,v_2\}$ and $\{v_1,v_3\}$.}
	\label{fig:many-bad}
	\end{center}
\end{figure}

The following proposition gives an example of a bipartite graph that admits many bad sets.

\begin{proposition}
	There exists a bipartite graph $G(U,V;E)$ with $n$ vertices on each side, which admits $2^{n/4}$ bad sets.
	\label{ex:many-bad}
\end{proposition}

\begin{proof}
	Consider the bipartite graph $G'$ (with four vertices on each side) depicted in \expref{Figure}{fig:many-bad}.
	The sets $V_1=\{v_1,v_3\}$ and $V_2=\{v_1,v_2\}$ are bad sets.
	Indeed, under $\pi=\{v_4,v_3,v_2,v_1\}$, the set $\{v_1,v_2\}$ remains unmatched if $u_1$ and $u_2$ arrive before $u_3$ and $u_4$.
	Similarly, under $\pi=\{v_4,v_2,v_3,v_1\}$, the set $\{v_1,v_3\}$ remains unmatched if $u_1$ and $u_2$ arrive before $u_3$ and $u_4$.
	Now, consider the graph $G$ that is composed of $n/4$ disjoint copies of $G'$.
	Since each copy admits two bad sets, $G$ admits $2^{n/4}$ bad sets.
\end{proof}


\section{Iterative process}
\label{sec:iterative}

A natural approach for establishing the existence of a good linear order $\pi$ is the following iterative process of ``upgrading'' unmatched vertices.
Given a linear order $\pi:V\rightarrow [n]$ and a linear order $\sigma: U\rightarrow [n]$, let $M_G[\sigma,\pi]$ be the result of the greedy matching where vertices in $U$ arrive in order $\sigma$ (from high to low) and each vertex $u\in U$ is matched to its highest (under $\pi$) neighbor (or left unmatched if all its neighbors are already matched).

\begin{figure}[h!]
	\begin{center}
\includegraphics[trim={0 0cm 0 0cm},clip,scale=.5]{figure5}
	\caption{An iterative process where unmatched vertices are given priority. In every iteration thick edges are in the matching; gray vertices are unmatched.}
	\label{fig:iterative-log}
	\end{center}
\end{figure}

Fix an arbitrary linear order $\pi_1$ on $V$, and let $\sigma_1$ be a linear order on $U$ minimizing the greedy matching.\footnote{It is unclear whether $\sigma_1$ can be computed in polynomial time. The related problem of computing a minimum maximal matching in bipartite graphs is known to be NP-hard \cite{ChlebikC06}. However, here we consider the existence problem.}  %
%
Let $M_1 = M_G[\sigma_1,\pi_1]$ be the result of the greedy matching under linear order $\sigma_1$ and $\pi_1$.
If $|M_1|/n$ is some constant greater than $1/2$, then terminate with linear order $\pi_1$.
Otherwise, partition $V$ into the set $V_{H_1}$ of unmatched vertices ($H$ for \emph{high}, as they will be placed high in the next iteration) and the set $V_{L_1}$ of matched vertices ($L$ for \emph{low}, as they will be placed low in the next iteration).

Consider now a linear order $\pi_2$ in which $V_{H_1}$ precedes $V_{L_1}$ (preserving the internal order between vertices in $V_{H_1}$ and similarly between vertices in $V_{L_1}$), and let $\sigma_2$ be a linear order on $U$ minimizing the resulting greedy matching.
Let $M_2=M_G[\sigma_2,\pi_2]$.
If $|M_2|/n$ is some constant greater than $1/2$, then terminate with linear order $\pi_2$.
Else, partition $V$ into the set $V_{H_2}$ of unmatched vertices and the set $V_{L_2}$ of matched vertices, and consider a linear order $\pi_3$ in which $V_{H_2}$ precedes $V_{L_2}$ (preserving internal orders).
Continue this iterative process until the
%
linear order $\pi_k$ obtained  %
ensures a matching greater than a half.

The intuition behind this approach is that the unmatched vertices need some ``help'' in order to be matched, and we provide this help in the form of
prioritizing them over matched vertices. %
%
%
One might hope that this process will reach a good linear order within a constant number of iterations.
Unfortunately, we show an example where the process goes through $\log n$ iterations before it first obtains a linear order ensuring a matching that exceeds $n/2$.

The construction of the graph is inductive.
The base is $G_0(U_0,V_0;E_0)$, with two vertices $u,v$ and a single edge between them.
For every $i=1,2,\ldots$, $G_i(U_i,V_i;E_i)$ is such that $|U_i|=|V_i|=2^i$; it is obtained by taking two (disjoint) copies of $G_{i-1}$, with additional edges of the form $(u_j,v_j)$ for every $u_j$ from one copy of $G_{i-1}$ to $v_j$ in the second copy of $G_{i-1}$.
An example of $G_3$ is presented in \expref{Figure}{fig:iterative-log}(a).
The iterative process is depicted in \expref{Figure}{fig:iterative-log}(a)-(d).
In all iterations preceding the last one, exactly $n/2$ vertices are matched in the worst $\sigma$.

 %
\section{Stable matching}
\label{sec:stable-matching}
The max-min greedy matching problem is also related to the stable matching problem~\cite{GS62} with \emph{partial preferences}.
In a stable matching scenario, every vertex in $U$ has a preference order over %
the vertices in $V$, and every vertex in $V$ has a preference order over %
the vertices in $U$. Given a matching $M$, a pair of vertices $u \in U, v \in V$ is said to constitute a blocking pair if $u$ and $v$ %
prefer each other over their partners in $M$. A matching is said to be stable if no pair of vertices $u,v$ constitutes a blocking pair. The stable matching problem is extended to partial preferences by including an option of being unmatched in the preference order, which can be preferred over being matched to some other vertex. This implies that if %
being  %
unmatched is preferred over being matched to $v$ by some vertex $u$, then $u$ being matched to $v$ cannot be stable.  
%
%
%
%

Consider a graph $G$ and linear orders $\pi$ and $\sigma$ over $V$ and $U$, respectively. In the corresponding stable matching with partial preferences instance, a vertex $v\in V$ sets its preferences over vertices in $N(v)\subseteq U$ according to the ranking in $\pi$, and prefers being unmatched over vertices in $U\setminus N(v)$. The analogous preferences are set for vertices in $U$ with respect to vertices in $V$ and $\sigma$. 
%
%
%
%
%
%
%
%
%
%
%
%
%
%
%
%
%
%
%
%
We show the following.

\begin{proposition}
	The partial preference orders defined by the graph $G$ and linear orders $\pi,\sigma$ imply the existence of a unique stable matching. 
	Moreover, the outcome of the greedy matching process is the unique stable matching. 
\end{proposition}
\begin{proof}
  First, we show that the preferences defined by $G$, $\pi$ and
%
$\sigma$   %
imply a unique stable matching. Let $\succ_\pi$ and $\succ_\sigma$ be the preference orders defined by $\pi$ and $\sigma$, where a higher ranking implies a higher preference. Let $u_1\succ_\sigma u_2\succ_\sigma\ldots\succ_\sigma u_n$ and $v_1\succ_\pi u_2\succ_\pi \ldots \succ_\pi v_n$ be the $U$-side and $V$-side vertices ordered according to $\succ_\sigma$ and $\succ_\pi$, respectively. Consider two \emph{different} matchings, $M_1$ and $M_2$ that are stable with respect to the preferences defined by $G$, $\pi$ and $\sigma$. Consider iterating over the $U$-side vertices according to their ordering until we reach a vertex $u_i$ with a different outcome in $M_1$ and $M_2$. There are two scenarios:
	\begin{itemize}
		\item $u_i$ is matched to $v_j$ in $M_1$ and to $v_k$ for $k\neq j$ in $M_2$. Assume without loss that $k<j$; that is, $v_k\succ_\pi v_j$ and also $v_k\succ_{u_i} v_j$. In this case we claim that $(u_i,v_k)$ is a blocking pair in $M_1$. To see this, we first notice $v_k$ is not matched to any $u_\ell$ with $\ell< i$ in $M_1$. Otherwise, $u_i$ would not have been the first $U$-side vertex for which $M_1$ and $M_2$ differ. Therefore, either $v_k$ is matched to some $u_\ell$ in $M_1$ for which $u_i\succ_\sigma u_\ell$ (and therefore $u_i\succ_{v_k} u_\ell$), or $v_k$ is unmatched in $M_1$. In both cases, since $v_k\succ_{u_i} v_j$, then both $u_i$ and $v_k$ \emph{strictly} improve their outcome in $M_1$ by matching to each other.
		\item $u_i$ is matched in one of the matchings and is unmatched in the other. Without loss of generality, $u_i$ is matched to $v_k$ in $M_2$ and is unmatched in $M_1$. As in the first scenario, since the matchings are identical until we reach $u_i$, either $v_k$ is unmatched in $M_1$, or it is matched to some $u_\ell$ for which $u_i\succ_{v_k} u_\ell$. This again implies that $(u_i,v_k)$ is a blocking pair in $M_1$ as both vertices improve their outcome by matching to each other.
	\end{itemize}
	
	To show that  the greedy matching process results in the unique stable matching, we argue there are no blocking pairs. Let $M$ be the resulting matching, and suppose towards a contradiction that $(u_i,v_j)$ is a blocking pair. This means that ($i$) the edge $(u_i,v_j)$ exists in $G$, and ($ii$) $v_j$ is not matched to a higher ranked vertex than $u_i$ in $M$. But this implies that when it is $u_i$'s turn in the greedy matching process, $v_j$ is an option. Therefore, it would not end up with an %
outcome, inferior  %
to being matched to $v_j$, a contradiction.
\end{proof}

%
Thus, the max-min greedy matching problem 
%
is equivalent to %
choosing a %
linear order $\pi$ over $V$ such that for every %
linear order $\sigma$ over $U$, the unique stable matching with respect to $\pi$ and $\sigma$ obtains a large matching.

%
%
%
%
%
%
%
%
%
%
%
%
%
%
%
%
%
%
%
%
%
%
%
%
%
%
%
%
%
%
%
%
%
%

%
%

\section{Additional results for regular graphs}
\label{sec:additional}

The following theorem shows that one cannot hope to get $\rho > 3/4$ with a random linear order in regular graphs.


%
\begin{theorem}
	\label{thm:3-4}
	For every $\epsilon > 0$ and sufficiently large $d$, there are $d$-regular graphs $G$ for which a random linear order $\pi$ results in $\rho \le \frac{3}{4} + \epsilon$.
\end{theorem}


\begin{proof}
	Consider a $d$-regular bipartite graph $G(U,V;E)$, where $U = U_1 \cup U_2$, $V = V_1 \cup V_2$,   $|U_1| = |V_1| = d$, $|U_2| = |V_2| = d-1$, and $|U| = |V| = 2d-1 = n$. The edge set is $(U_1 \times V_2) \cup (U_2 \times V_1) \cup M$, where $M$ is a perfect matching between $U_1$ and $V_1$. Hence the degree of every vertex is $d$, and $(U_2,V_2)$ forms 
	a balanced bipartite independent set.
	
	Let $Q_1 \subset V$ (a random variable) be the set of first $d$ vertices under the random linear order $\pi$, and let $Q_2 = V \setminus Q_1$. Then, $E[|V_2 \cap Q_2|] =  (\frac{d-1}{n})^2 n > \frac{n}{4} - 1$. Moreover, as the value of $|V_2 \cap Q_2|$ is concentrated around its mean, then for sufficiently large $d$ (say,
%
$d = \Omega\left(1/\epsilon^2)\right)$),    %
the actual value of $|V_2 \cap Q_2|$ is above 
%
$(1 - \epsilon)n/4$  %
with high probability. Hence the bounds we prove hold with high probability, and not just in expectation.
	
	Observe that there is a perfect matching $M'$ between $U_1$ and $Q_1$: For every vertex $u\in U_1$, if its matched vertex under $M$ is in $V_1 \cap Q_1$ then match $u$ to this vertex, and otherwise match $u$ to an arbitrary vertex in $V_2 \cap Q_1$. The number of vertices in $U_1$ whose partner under $M$ is not in $Q_1$ equals $|V_2 \cap Q_1|$, and so this gives a perfect matching. 
	For $1 \le j \le d$, let $u_j$ denote the vertex in $U_1$ that is matched to the $j$th vertex in $Q_1$ under $M'$ according to $\pi$. 
	Choose a linear order $\sigma$ over $U$ that begins with $u_1, \ldots, u_d$, and is followed by the vertices of $U_2$ (unlike the vertices of $U_1$, the vertices of $U_2$ can be arranged in an arbitrary order). With this $\sigma$, all vertices of $U_1$ will be matched with $Q_1$, and consequently no vertex of $V_2 \cap Q_2$ will be matched. This leaves
%
$(1 - \epsilon)n/4$   %
vertices unmatched, proving the theorem.
\end{proof}


%
%
%
%

%
%
%
%

%
%
%
%
%
%
%

We also establish a few impossibility results for regular graphs of low degree.

\begin{theorem}
	The following hold.  %
	\begin{itemize}
		\item There exists a $3$-regular bipartite graph $G$ for which $\rho(G) = \frac{5}{7}$.
		\item There exists a $4$-regular bipartite graph $G$ for which $\rho(G) = \frac{10}{13}$.
	\end{itemize}
	\label{thm:imposs-projective}
\end{theorem}

%
%
%
The proof relies on %
the incidence graphs of finite  %
projective planes.
A projective plane consists of a set of lines and a set of points, where
%
every two lines intersect in a single point, %
every two points are incident to a single line,
and there are four points, no three of which are on a line.  %
A projective plane %
defines %
a bipartite graph $G(U,V;E)$,
called its \emph{incidence graph},  %
where every vertex $u \in U$ corresponds to a point in the plane, every vertex $v \in V$ corresponds to a line, and there exists an edge between $u$ and $v$ if the point corresponding to $u$ is incident to the line corresponding to $v$.

\begin{proof}
  For the first result, we show that $\rho = \frac{5}{7}$ for
%
incidence graph of   %
the Fano plane.
	The Fano plane is a projective plane consisting of 7 points and 7 lines, with 3 points on every line and 3 lines through every point.
%
%
Let $G(U,V;E)$ be the incidence graph of the Fano plane.  This is a        
$3$-regular bipartite graph.   %
Let $N(V')$ denote the
set of  %
neighbors of a set $V' \subset V$.
For every set $V' \subset V$ %
%
of size $|V'|=2$
%
 we have $|N(V')|=5$.   %
We show below that for every such $V'$ there exists a perfect matching between $N(V')$ and $V \setminus V'$.
	Hence one can choose a linear order $\sigma$ over $U$ whose first $5$ vertices are $N(V')$ that will match the vertices of $V \setminus V'$ one by one. Thereafter, the vertices of $V'$ will remain unmatched.
By Hall's condition, it suffices to show that for every set $U' \subset N(V')$
%
of size %
$|U'| \leq 5$ %
we have  %
$|N(U')| \geq |U'|+2$ (so that Hall's condition %
holds %
with respect to the set $V \setminus V'$).
	Indeed, for every set $U'$ of size 1, $|N(U')|=3$, for every set $U'$ of size %
2, $|N(U')|=5$, for every set $U'$ of size $\geq 3$, $|N(U')|\ge 6$, %
and for every set $U'$ of size $5$, $|N(U')|=7$.
%
These statements can be quickly verified using the symmetries
of the Fano plane.  %
	It follows that $\rho(G)=5/7$.
	
The second result follows %
%
by  %
a similar argument.
	It is known that there exists a projective plane consisting of $13$ points and $13$ lines, with $4$ points on every line and $4$ lines through every point.
	We claim that $\rho = \frac{10}{13}$ for the
%
incidence graph $G(U,V;E)$ of %
this projective plane.
	By the properties of a projective plane, for every set 
%
$V' \subset V$   %
such that $|V'|=3$, it holds that $|N(V')| \in \{9,10\}$.
	We show below that for every such $V'$ there exists a perfect matching between $N(V')$ (and possibly an additional vertex $u$ in case $|N(V')| = 9$) and $V \setminus V'$.
	Hence one can choose a linear order $\sigma$ over $U$ whose first $10$ vertices are $N(V')$ (possibly with the additional vertex) that will match the vertices of $V \setminus V'$ one by one. Thereafter, the vertices of $V'$ will remain unmatched.
	By Hall's condition, it suffices to show that for every set $U' \subset N(V')$ such that $|U'| \leq 10$ it holds that $|N(U')| \geq |U'|+3$ (so that Hall's condition %
holds  %
with respect to the set $V \setminus V'$).
	Indeed, for every set $U'$ of size 1, $|N(U')|=4$, for every set $U'$ of size $\geq 2$, $|N(U')| \geq 7$, for every set $U'$ of size $\geq 5$, $|N(U')| \geq 11$,
%
%
%
for every set $U'$ of size $9$, $|N(U')| \geq 12$,
and for every set $U'$ of size $\geq 10$, $|N(U')| = 13$.  %
	It follows that $\rho(G)=10/13$.
\end{proof}

\section{Open problems}
	In this paper we proved various bounds regarding the max-min greedy matching problem. Our main result was in showing that for every bipartite graph $G$, one can efficiently compute a linear order $\pi$ such that for every $\sigma$, $\rho$ is %
greater than  %
 a constant strictly greater than $1/2$ (namely, greater than $0.51$). In~\cite{Cohen-AddadEFF16}, they present a $G$ such that for every $\pi$, there is a $\sigma$ such that $\rho \leq 2/3$. The main open question left by this paper is to close the gap between $0.51$ and $2/3$. Moreover, is it possible to compute the optimal $\pi$ efficiently? 
	
 In addition to our result for general bipartite graphs, we also have results for interesting families of bipartite graphs. Specifically, we show that for regular %
%
graphs and for  %
Hamiltonian graphs, we are able to compute linear orders with improved bounds. Namely, we show that for regular graphs, $\rho\geq 5/9-O(1/\sqrt{d})$, and for Hamiltonian graphs, $\rho \ge 5/9$. Moreover, as opposed to general graphs, for these families of graphs, a random linear order guarantees 
that $\rho$ is %
greater than %
a constant strictly %
greater  %
than $1/2$ (by some tiny quantity). However, there are substantial gaps remaining for these families of graphs as well. The $2/3$ example from~\cite{Cohen-AddadEFF16} is of a Hamiltonian graph, while the only upper bound known on the value of $\rho$ for regular graphs is $8/9$. It will be interesting to see what
are the best achievable bounds on $\rho$,
and what are the best guarantees that a random $\pi$ can achieve. 
	
Finally, the results in our paper can be interpreted as welfare bounds that a seller can achieve by pricing items in full information markets with binary unit-demand valuations. One can ask the same question for more general classes of valuations; for instance, general unit-demand valuations. In this more general version of the problem, the bipartite graph is now weighted, and the seller, knowing the graph, sets prices over the right hand side. Then, the left hand side arrives in some order, and each arriving ``buyer'' on the left hand side chooses an available ``item'' on the right hand side that maximizes the weight of the edge \emph{minus} the price of the item. For this more general problem, it is still unknown if the seller can set prices that guarantee that the weight of the resulting weighted matching is at least a $\rho$-fraction of the maximum weight matching for some constant $\rho>1/2$.
%
%
%
In fact, $2/3$ might actually be the correct answer for this
more general problem.   %


\section*{Acknowledgments}

	A substantial part of this work was conducted in Microsoft Research. We are grateful to Amos Fiat and Sella Nevo for numerous discussions that contributed significantly to the ideas presented in this paper. We also thank Robert Kleinberg for helpful discussions.

\nocite{conf-version}
